

\section{Proofs} \label{sec-proofs-high-d}

We will prove Theorem \ref{thm:bipartite-random} using the same general strategy as in \cite{HLV2015}.  Specifically, the proof of Theorem \ref{thm:bipartite-random} can be separated into two parts:  a recovery guarantee under a set of deterministic conditions, and a proof that the random model meets these conditions with high probability.  These sufficient deterministic conditions, roughly speaking, are (1) that the graph is connected and the nodes have tightly controlled degrees; (2) that the camera and structure locations are all distinct;  (3) that all pairwise distances between cameras and locations are within a constant factor of each other; (4) that any choice of two camera locations and two structure locations live in a three dimensional affine space; (5) that the camera and structure locations are `well-distributed' in a sense that we will make precise; and (6) that there are not too many corruptions affecting a single camera location or structure point.  Theorem \ref{thm:mainthm} in Section \ref{sec-deterministic-theorem} states these deterministic conditions formally.

As in \cite{HLV2015}, we will prove the deterministic recovery theorem directly, using several geometric properties concerning how deformations of a set of points induce rotations. 
%Note that an infinitesimal rigid rotation of two points $\{t_i, t_j\}$ about their midpoint to ${\{t_i + h_i, t_j + h_j\}}$ is such that $h_i - h_j$ is orthogonal to $t_{ij} = t_i - t_j$.  We will abuse terminology and say that $\|P_{\tij^\perp} (h_j - h_i)\|$ is a measure of the rotation of a finite deformation $\{h_i, h_j\}$. Morally speaking, because ShapeFit attempts to minimizes the total rotations within a feasible deformation, no deformation of the true locations results in an equal or lower objective.  These geometric properties are as follows: 
%\red{[Would like to rearrange the whole paragraph as follows:]} 
Note that an infinitesimal rigid rotation of two points $\{t_i, t_j\}$ about their midpoint to ${\{t_i + h_i, t_j + h_j\}}$ is such that $h_i - h_j$ is orthogonal to $t_{ij} = t_i - t_j$.  We will abuse terminology and say that $\|P_{\tij^\perp} (h_i - h_j)\|$ is a measure of the rotation in a finite deformation $\{h_i, h_j\}$, and we say that $\langle h_i - h_j, \ti - \tj \rangle$ is the amount of stretching in that deformation. Using this terminology, the geometric properties we establish are:

\begin{itemize}
\item If a deformation stretches two adjacent sides of a $C_4$ at different rates, then that induces a rotation in some edge %the rest 
of the $C_4$ (Lemma \ref{lem:rigidity4}).  
\item If a deformation rotates one edge shared by many $C_4$s, then it induces a rotation over many of those $C_4$s, provided the opposite points of those triangles are `well-distributed' (Lemma \ref{lem:c4s}).
\item A deformation that rotates bad edges, must also rotate good edges (Lemma \ref{lem:badgood-bipartite}).
\item For any deformation, some fraction of the sum of all rotations must  affect the good edges (Lemma \ref{lem:transference-bipartite}).
\end{itemize}
By using these geometric properties, we show that all nonzero feasible deformations induce a
large amount of total rotation. Since some fraction of the total rotation must be on the good edges,  the objective must increase. 

The main technical difference between the present proof and the proof of \cite{HLV2015} is that the proof in $\cite{HLV2015}$ relies on the presence of many triangles in the graph of uncorrupted measurements.  Because of the bipartite structure of the present work, there are no triangles in the graph.  Hence, the technical novelty of the present proof is the establishment of the properties above when there are a sufficient number of $C_4$s in the graph of uncorrupted measurements.  

In Section \ref{sec-deterministic-theorem}, we present the deterministic recovery theorem.  In Section \ref{sec-unbalanced-motions}, we present and prove Lemma \ref{lem:rigidity4}.  In Section \ref{sec-c4-inequality}, we present and prove Lemmas \ref{lem:c4s}--\ref{lem:transference-bipartite}.  In Section \ref{sec-proof-mainthm}, we prove the deterministic recovery theorem.  In Section \ref{sec-properties-of-gaussians}, we prove that Gaussians satisfy several properties with high probability. In Section \ref{sec-gaussians-well-distributed}, we prove that Gaussians satisfy well-distributedness with high probability.  In Section \ref{sec-random-graph}, we prove that \erdosrenyi graphs are connected and have controlled degrees and codegrees with high probability.  Finally, in Section \ref{sec-proof-random-theorem}, we prove Theorem \ref{thm:bipartite-random}.


%%
%%
%%
%%
%%
%% Deterministic recovery condition
%%
%%
%%
%%
%%
%%


\subsection{Deterministic recovery theorem in high dimensions} \label{sec-deterministic-theorem}

To state the deterministic recovery theorem, we need two definitions.
%Let $|\Vl| = \nl$ and $|\Vs|=\ns$.  Let $\Knn $ be the bipartite-complete graph over $\nl + \ns$ vertices.
The first definition captures the `regularity' of the measurement graph. A random bipartite graph can easily be seen to satisfy the conditions. Note that the definition does not depend on the vectors locations $\{t_i\}$ and $\{p_j\}$.

\begin{definition}
We say that a graph $G(\Vl \cup \Vs,E)$ is \emph{bipartite-$p$-typical} if it satisfies the following
properties:
\begin{enumerate}
  \setlength{\itemsep}{1pt} \setlength{\parskip}{0pt}
  \setlength{\parsep}{0pt}
	\item $G$ is connected,
	\item each vertex in $\Vl$ has degree between $\frac{1}{2}\ns p$ and $2\ns p$, and\\ each vertex in $\Vs$ has degree between $\frac{1}{2}\nl p$ and $2 \nl p$.
	\item each pair of vertices in $\Vl$ has codegree between $\frac{1}{2}\ns p^2$ and $2 \ns p^2$, where the codegree  of $j,l \in \Vl = |\{i \mid ij \in E(G), il \in E(G) \} |$.  Each pair of vertices in $\Vs$ has codegree between $\frac{1}{2}\nl p^2$ and $2 \nl p^2$.
\end{enumerate}
\end{definition}
%Note that if $G$ is $p$-typical, then its number of edges is between $\frac{1}{4}n^{2}p$ and $n^{2}p$.


The next definition captures how `well-distributed' the location points $\{t_i\}$ and $\{p_j\}$ are in $\R^d$.  

\begin{definition} \ 
\begin{itemize}
  \setlength{\itemsep}{1pt} \setlength{\parskip}{0pt}
  \setlength{\parsep}{0pt}
\item[(i)] Let $S = \{(t_k, p_k)\}_{k=1\ldots m} \subset \R^d \times \R^d$. Let $x,y \in \R^d$.
We say that $S$ is $c$-well-distributed with respect to $(x,y)$ if the following holds for all $h \in \R^d$:
\[
	\sum_{(t,p) \in S}\|\proj_{\Span\{p-x,t-p, y-t\}^{\perp}}(h)\|_2 \ge c|S|\cdot\|\proj_{(x-y)^{\perp}}(h)\|_2.
\]
\item[(ii)] Let $T = \{\ti\}_{i \in \Vl}$ and $P = \{\pj\}_{j \in \Vs}$.  We say that $(T,P)$ is c-well-distributed along $G$ if for all $i \in \Vl, j \in \Vs$, the set $S_{ij} = \{(t_k, p_\ell)\,:\, i\ell \in E(G), k\ell \in E(G), kj \in E(G), k \neq i, \ell \neq j \}$ is $c$-well-distributed with respect to $(t_i, p_j)$.
\end{itemize}
\end{definition}


We now state sufficient deterministic recovery conditions on the graph $G$, the subgraph $\Eb$ corresponding to corrupted observations, and the locations $\Tnot$ and $\Pnot$.

\begin{theorem} \label{thm:mainthm}
Suppose $\Tnot, \Pnot, \Eb, G$ satisfy the conditions
\begin{enumerate}
  \setlength{\itemsep}{1pt} \setlength{\parskip}{0pt}
  \setlength{\parsep}{0pt}
	\item The underlying graph $G$ is bipartite-$p$-typical,
	\item All vectors in $\Tnot$, $\Pnot$ and $\Tnot \cup \Pnot$ are distinct, respectively.
	\item For all $i,k \in \Vl$ and $j, \ell \in \Vs$, we have $c_{0}\|t_{k\ell}^{(0)}\|_2\le\|t_{ij}^{(0)}\|_2$,
	\item For all $i,k \in \Vl, j,\ell \in \Vs$ such that $k \neq i, j \neq \ell$, we have \\
	$\min \Bigl(\|\proj_{\Span (\tokj,\toil)^{\perp}} \toij \|_2, 
	 \|\proj_{\Span (\tokl,\toil)^{\perp}} \toij \|_2 \Bigr) / \|\toij \|_2 \geq \beta$
%	 $\min \Bigl(\|\proj_{\Span (\toik,\tojl)^{\perp}} \toij \|_2, 
%	 \|\proj_{\Span (\tojl,\tokl)^{\perp}} \toij \|_2 \Bigr) / \|\toij \|_2 \geq \beta$
	\item The pair $(\Tnot, \Pnot)$ is $c_{1}$-well-distributed
along $G$,
	\item Each vertex in $\Vl$ (resp.  $\Vs$) has at most $\eps \ns$ (resp.  $\eps \nl$) incident edges in $\Eb$.
\end{enumerate}
for constants $0< p, c_0, \beta,  c_{1}, \eps \leq 1$.  If $\eps \leq \frac{\beta c_0 c_1^2 p^4}{384\cdot204\cdot64}$ and $\nl, \ns > max(64, \frac{8}{p^2})$, then $L(\Tnot, \Pnot)\neq 0$ and $(\Tnot, \Pnot) / L(\Tnot, \Pnot)$ is the unique optimizer of ShapeFit.
\end{theorem}

%\red{A cycle of length $4$, denoted as $C_4$, will be the main substructure used in proving Theorem~\ref{thm:mainthm}.
%It is well-known that $C_4$ is parallely-rigid \blue{((we don't define parallel rigid))} in $\R^d$ for $d \ge 3$ unless the given vectors
%are `degenerate'. The parameter $\beta$ in Condition 4 can be understood as a measure of how far the 4-tuples of vectors $t_{ij}^{(0)}, t_{kj}^{(0)}, t_{i\ell}^{(0)}, t_{k\ell}^{(0)}$
%for $i,k \in \Vl$, $j,\ell \in \Vs$ are from being `degenerate' in this sense (we will make this
%precise in Subsection \ref{sec-unbalanced-motions}). }
Before we prove the theorem, we establish that $L(\Tnot, \Pnot)\neq 0$ when $\eps$ is small enough.  This property guarantees that some scaling of $(\Tnot,\Pnot)$ is feasible and occurs, roughly speaking, when $|\Eb| < |\Eg|$.


\begin{lemma} \label{lem:T0-feasible-bipartite} 
%\red{If $|E_b| < c_0 |\Eg|$, then $L(\Tnot) \neq 0$.  Further, if $\varepsilon < \frac{c_0 p}{8}$, then $|E_b| < c_0 |\Eg|$.}
If $\varepsilon < \frac{c_0 p}{4}$, then $L(\Tnot, \Pnot) \neq 0$.
\end{lemma}
\begin{proof}
Since $v_{ij} = \toijhat$ for all $ij \in E_g$, we have
\[
L(\Tnot,\Pnot) = \sum_{ij \in E(G)} \langle \toij, v_{ij} \rangle \geq \sum_{ij \in E_g} \| \toij\|_2 - \sum_{ij \in E_b} \| \toij \|_2.
\]
By Condition 3,  $c_0 \muinf |E_g| \leq \sum_{ij \in E_g} \| \toij\|_2$ and $\muinf |E_b| \geq  \sum_{ij\in E_b} \| \toij\|_2$.  Thus it suffices to prove that $c_0 |E_g| > |E_b|$. 
As $\varepsilon < \frac{p}{4}$, Condition 1 and 6 gives $|E_g| \geq \frac{1}{2} \nl \ns p - \eps \nl \ns \geq \frac{1}{4}\nl \ns p$. Since $|E_b| \leq \eps \nl \ns$, if $\eps < \frac{c_0 p}{4}$, then we have $c_0 |E_g| > |E_b|$. 
\end{proof}




\subsection{Unbalanced parallel motions induce rotation} \label{sec-unbalanced-motions}

The following lemma concerns geometric properties of deformations of a set of points.  Specifically it shows that if four points are deformed in a way that differentially scales the lengths of two edges, then it necessarily induces a rotation somewhere in a $C_4$ containing those points.

%\red{As briefly mentioned in the previous subsection, $C_4$ is known to be parallel-rigid \blue{((we don't define parallel rigid))} in
%$\R^d$ for $d \ge 3$. In other words, given four vectors $t_{12}, t_{23}, t_{34}, t_{41} \in \R^d$ (for $d \ge 3$), there is a unique $4$-tuple of vectors $t_1,t_2,t_3,t_4 \in \R^d$ satisfying $t_{i(i+1)} = t_i - t_{i+1}$ for all $i \in [4]$ up to scale and translation, unless the given vectors are `degenerate' (index summation is modulo 4). The following lemma can be understood as a tool measuring the distance of the given $4$-tuple of vectors from being `degenerate'.}

\begin{lemma} \label{lem:rigidity4}
Let $d\ge3$. Let $t_{1},t_{2},t_{3},t_{4} \in \R^d$ be distinct.  Let $t_{ij} = \ti - \tj$ and $\tijhat = \frac{\tij}{\| \tij \|}$.  Let $v_{1},v_{2},v_{3},v_{4}\in\mathbb{R}^{d}$ and $\alpha \in \R$.   Let $\{\deltatilde_{i(i+1)}\}$ be such that  ${\langle v_{i}-v_{i+1}-\alpha t_{i(i+1)},\hat{t}_{i(i+1)}\rangle=}$ $\deltatilde_{i(i+1)}\|t_{i(i+1)}\|_2$ for each $i\in[4]$, where index summation is modulo 4.
\begin{align*}
(i) \quad &\sum_{i\in[4]} 
	\|\proj_{t_{i(i+1)}^{\perp}}(v_{i}-v_{i+1})\|_2
	\ge \left\| \proj_{\Span(t_{23},t_{41})^{\perp}}t_{12}\right\|_2 \left|\deltatilde_{12}-\deltatilde_{34}\right| \\
(ii) \quad &\sum_{i\in[4]} 
	\|\proj_{t_{i(i+1)}^{\perp}}(v_{i}-v_{i+1})\|_2
	\ge  \left\| \proj_{\Span(t_{34},t_{41})^{\perp}} t_{12}\right\|_2 \left|\deltatilde_{12}-\deltatilde_{23}\right|.
\end{align*}
\end{lemma}
\begin{proof}
%Note that $t_{ij}=-t_{ji}$ and $\deltatilde_{ij}=\deltatilde_{ji}$ for each distinct $i,j\in[4]$. 
The given condition implies $\proj_{t_{i(i+1)}^{\perp}}(v_{i}-v_{i+1})=v_{i}-v_{i+1}-(\alpha+\deltatilde_{i(i+1)})t_{i(i+1)}$
for each $i \in [4]$. Therefore,
\begin{eqnarray}
\sum_{i \in [4]}\|\proj_{t_{i(i+1)}^{\perp}}(v_{i}-v_{i+1})\|_2 & = & \sum_{i \in [4]}\left\| v_{i}-v_{i+1}-\left(\alpha+\deltatilde_{i(i+1)}\right)t_{i(i+1)}\right\|_2 \nonumber \\
 & \ge & \left\| \sum_{i \in [4] }v_{i}-v_{i+1}-\left(\alpha+\deltatilde_{i(i+1)}\right)t_{i(i+1)}\right\|_2 \nonumber \\
 & = & \|\deltatilde_{12}t_{12}+\deltatilde_{23}t_{23}+\deltatilde_{34}t_{34}+\deltatilde_{41}t_{41}\|_2. \label{eq:rigidity}
\end{eqnarray}

\noindent (i) Since $\deltatilde_{34}(t_{12}+t_{23}+t_{34}+t_{41})=0$, the right-hand-side of
\eqref{eq:rigidity} equals $\|(\deltatilde_{12}-\deltatilde_{34})t_{12}+(\deltatilde_{23}-\deltatilde_{34})t_{23}+(\deltatilde_{41}-\deltatilde_{34})t_{41}\|_2$.
The conclusion follows since
\begin{eqnarray*}
\left\| (\deltatilde_{12}-\deltatilde_{34})t_{12}+(\deltatilde_{23}-\deltatilde_{34})t_{23}+(\deltatilde_{41}-\deltatilde_{34})t_{41}\right\|_2  & \ge & \min_{s,s'\in\mathbb{R}}\|(\deltatilde_{12}-\deltatilde_{34})t_{12}-st_{23}-s't_{41}\|_2\\
 & = & \left\| \proj_{\Span(t_{23},t_{41})^{\perp}}(\deltatilde_{12}-\deltatilde_{34})t_{12}\right\|_2.
\end{eqnarray*}

\noindent (ii) Since $\deltatilde_{23}(t_{12}+t_{23}+t_{34}+t_{41})=0$, the right-hand-side
of \eqref{eq:rigidity} equals $\|(\deltatilde_{12}-\deltatilde_{23})t_{12}+(\deltatilde_{34}-\deltatilde_{23})t_{34}+(\deltatilde_{41}-\deltatilde_{23})t_{41}\|_2$.
The conclusion follows since
\begin{eqnarray*}
\left\| (\deltatilde_{12}-\deltatilde_{23})t_{12}+(\deltatilde_{34}-\deltatilde_{23})t_{34}+(\deltatilde_{41}-\deltatilde_{23})t_{41}\right\|_2  & \ge & \min_{s,s'\in\mathbb{R}}\|(\deltatilde_{12}-\deltatilde_{23})t_{12}-st_{34}-s't_{41}\|_2\\
 & = & \left\| \proj_{\Span(t_{34},t_{41})^{\perp}}(\deltatilde_{12}-\deltatilde_{23})t_{12}\right\|_2.\qedhere
\end{eqnarray*}
\end{proof}



\subsection{$C_4$s inequality and rotation propagation} \label{sec-c4-inequality}

The following lemma is a generalization of the triangle inequality in a context of the rotational part of structure deformations.
%\red{The following lemma can be interpreted as a generalization of the triangle inequality, since
%the conclusion immediately follows from the triangle inequality if we remove the projectors in both sides of the inequalities. }

\begin{lemma}[$C_4$s Inequality] \label{lem:c4s}
Let $d\ge4$; $x,y \in \mathbb{R}^{d}$.  Let $S=\{(t_1, p_1), \cdots, (t_k, p_k)\} \subset \R^d \times \R^d$.  If $S$  is $c$-well-distributed
with respect to $(x,y)$, then for all vectors $h_{x},h_{y},h_{t_1},\cdots,h_{t_k},h_{p_1},\cdots,h_{p_k}\in\mathbb{R}^{d}$
and sets $X\subseteq[k]$, we have
\[
\sum_{i\in[k]\setminus X}\|\proj_{(x-p_{i})^{\perp}}(h_{x}-h_{p_i})\|_2 + 
\|\proj_{(p_{i}-t_i)^{\perp}}(h_{p_i}-h_{t_i})\|_2 + \|\proj_{(t_{i}-y)^{\perp}}(h_{t_i}-h_{y})\|_2 \ge(ck-|X|)\cdot\|\proj_{(x-y)^{\perp}}(h_{x}-h_{y})\|_2.
\]
\end{lemma}
\begin{proof}
For each $i\in[k]$, define $W_i = \Span(x - p_i, p_i - t_i, t_i - y)$.
Define $P$ as the projection map to the space of vectors orthogonal
to $x-y$, and define $P_{i}$ for each $i\in[k]$  as the projection
map to $W_{i}^{\perp}$. Since $(x-p_{i})^{\perp}\supseteq W_{i}^{\perp}$, $(p_i - t_i)^{\perp}\supseteq W_{i}^{\perp}$, and $(t_i - y)^{\perp}\supseteq W_{i}^{\perp}$,
it follows that 
\begin{eqnarray*}
 &  &\sum_{i\in[k]\setminus X}\|\proj_{(x-p_{i})^{\perp}}(h_{x}-h_{p_i})\|_2 + 
\|\proj_{(p_{i}-t_i)^{\perp}}(h_{p_i}-h_{t_i})\|_2 + \|\proj_{(t_{i}-y)^{\perp}}(h_{t_i}-h_{y})\|_2 \\
 & \ge & \sum_{i\in[k]\setminus X}\|P_{i}(h_{x}-h_{p_i})\|_2 +\|P_{i}(h_{p_i}-h_{t_i})\|_2 +\|P_{i}(h_{t_i}-h_{y})\|_2 \ge\sum_{i\in[k]\setminus X}\|P_{i}(h_{x}-h_{y})\|_2.
\end{eqnarray*}
Since $\{(t_1, p_1), \cdots, (t_k, p_k)\}$ are well-distributed with respect to $(x,y)$,
we have
\begin{equation}
\sum_{i\in[k]}\|P_{i}(h_{x}-h_{y})\|_2 \ge ck\cdot\|P(h_{x}-h_{y})\|_2.\label{eq:proj_triangles}
\end{equation}
Since $(x-y)^\perp \supseteq W_i^\perp$, we have $\|P_{i}(h_{x}-h_{y})\|_2\le\|P(h_{x}-h_{y})\|_2$ for all $i$. Hence
\[
\sum_{i\in[k]\setminus X}\|P_{i}(h_{x}-h_{y})\|_2 \ge(ck-|X|)\cdot\|P(h_{x}-h_{y})\|_2,
\]
proving the lemma.
\end{proof}

The proof of Theorem \ref{thm:mainthm} will rely on the following two lemmas, which state that rotational motions on some parts of the graph bound rotational motions on other parts. The following lemma relates the rotational motions on bad edges to the rotational motions on good edges.  Recall the notation $\tij = (1 + \deltaij) \toij + \etaij \sij$ where $\sij$ is a unit vector orthogonal to $\toij$ and $\etaij = \| P_{\toijperp} \tij \|_2$. 



\begin{lemma} \label{lem:badgood-bipartite} 
Fix $T, P$. If $\eps \leq \frac{c_1 p^3}{48}$ and $p \geq \sqrt{\frac{8}{n}}$, then $\sum_{ij\in E_{g}}\eta_{ij}\ge\frac{c_{1}p^{3}}{48\eps}\sum_{ij\in E_{b}}\eta_{ij}$. 
\end{lemma}
\begin{proof}
Let $i \in \Vl, j \in \Vs$.  Note that Condition 1 implies $|\{(k,\ell) \mid k\neq i; \ell \neq j; i\ell,k\ell,kj\in E(G) \}|\geq (\frac{1}{2} \ns p -1) (\frac{1}{2} \nl p^2 - 1) \geq \frac{1}{8} \nl \ns p^3$ if $p \geq \sqrt{\frac{8}{n}}$. By Condition 6, the number of pairs $(k,\ell) \in \Vl \times \Vs$ such that at least one of the edges $i\ell, k\ell, kj$ are in $E_b$ can be counted by considering the case when $i\ell \in E_b$ (at most $(\varepsilon \ns) \nl$ pairs), $kj \in E_b$ (at most $(\varepsilon \nl) \ns$ pairs), and $k\ell \in E_b$ (at most $\varepsilon \ns \nl$ pairs). Hence in total, there are at most $3\varepsilon \ns \nl$ such pairs.
  By Lemma \ref{lem:c4s}, the $c_1$-well-distributedness of $(\Tnot,\Pnot)$ along $G$, and the assumption that $\eps \leq  \frac{c_1 p^3}{48}$, we have 
\[
	\sum_{\substack{k \in \Vl, \ell \in \Vs\\k\neq i,\ell \neq j\\i\ell, k\ell, kj\in E_{g}}}(\eta_{i\ell}+\eta_{k\ell} + \eta_{kj})
	\ge\left(c_{1}\cdot\frac{1}{8}\nl \ns p^3- 3 \eps \nl \ns\right)\cdot\eta_{ij}
	\ge\frac{c_{1}}{16}\nl \ns p^3\cdot\eta_{ij}.
\]
Therefore, if we sum the inequality above for all bad edges $ij\in E_{b}$, then
\[
	\sum_{ij\in E_{b}}
	\sum_{\substack{k \in \Vl, \ell \in \Vs\\k\neq i,\ell \neq j\\il, k\ell, kj\in E_{g}}}(\eta_{i\ell}+\eta_{k\ell} + \eta_{kj})
\ge\frac{c_{1}}{16}\nl \ns p^3\cdot\sum_{ij\in E_{b}}\eta_{ij}.
\]
%Note that by condition 4, we have that for any $i\in\Vl$, the number of adjacent bad edges is at most $\eps \ns$, and we have that for any $j \in \Vs$, the number of adjacent bad edges is at most $\eps \nl$.  Hence, $|\Eb| \leq \eps \nl \ns$. 

For fixed $k\ell\in E_{g}$, the left-hand-side may sum $\eta_{k\ell}$
as many times as the number of $C_4$s in $E(G)$ that contain $k\ell$ and exactly one bad edge.  This is the same as the number of $C_4$s whose edge opposite $k\ell$ is bad, plus the number of $C_4$s whose edge adjacent to $\ell$ is bad, plus the number of $C_4$s whose edge adjacent to $k$ is bad.   In each case, there are at most $\eps \nl \ns$ such $C_4$s.
Hence, the left-hand-side of above is at most 
\[
	\sum_{ij\in E_{b}}
	\sum_{\substack{k \in \Vl, \ell \in \Vs\\k\neq i,\ell \neq j\\i\ell, k\ell, kj\in E_{g}}}(\eta_{i\ell}+\eta_{k\ell} + \eta_{kj})
	\le 3 \eps \nl \ns\cdot\sum_{ij\in E_{g}}\eta_{ij}.
\]
Therefore by combining the two inequalities above, we obtain 
\[
	\sum_{ij\in E_{b}}\eta_{ij}
	\le\frac{48\eps}{c_{1}p^{3}}\sum_{ij\in E_{g}}\eta_{ij}. %\leq \sum_{ij\in E_{g}}\eta_{ij}.
	\qedhere
\]
\end{proof}

The following lemma relates the rotational motions over the good graph $\Eg$ to rotational motions over the complete bipartite graph $\Knn$.

\begin{lemma} \label{lem:transference-bipartite} 
Fix $T, P$. If $\eps \leq \frac{c_1 p^3}{48}$ and $p \geq \sqrt{\frac{8}{n}}$, then $\sum_{ij\in E_{g}}\eta_{ij}\ge\frac{c_{1}p}{192}\sum_{ij\in E(\Knn)}\eta_{ij}$.
\end{lemma}
\begin{proof}
%\red{Let $K = \Knn$.}
Let $i \in \Vl, j \in \Vs$.  Note that Condition 1 implies $|\{(k,\ell) \mid k\neq i; \ell \neq j; i\ell,k\ell,kj\in E(G) \}|\geq (\frac{1}{2} \ns p -1) (\frac{1}{2} \nl p^2 - 1) \geq \frac{1}{8} \nl \ns p^3$ if $p \geq \sqrt{\frac{8}{n}}$.  Similarly as in Lemma~\ref{lem:badgood-bipartite}, Condition 6 implies that
the number of pairs $(k,\ell) \in \Vl \times \Vs$ such that at least one of the edges $i\ell, k\ell, kj$ are in $E_b$ is at most $3\eps \nl \ns$.
%By Condition 6, $|\Eb| \leq \eps \nl \ns$.  
By Lemma \ref{lem:c4s}, the $c_1$-well-distributedness of $(\Tnot,\Pnot)$ along $G$, and the assumption that $\eps \leq  \frac{c_1 p^3}{48}$, we have 
\[
	\sum_{\substack{k \in \Vl, \ell \in \Vs\\k\neq i,\ell \neq j\\i\ell, k\ell, kj\in E_{g}}}(\eta_{i\ell}+\eta_{k\ell} + \eta_{kj})
	\ge\left(c_{1}\cdot\frac{1}{8}\nl \ns p^3- 3 \eps \nl \ns\right)\cdot\eta_{ij}
	\ge\frac{c_{1}}{16}\nl \ns p^3\cdot\eta_{ij}.
\]
Therefore, if we sum the inequality above for all $i\in \Vl, j \in \Vs$, or equivalently over all $ij \in E(\Knn)$, then
\[
	\sum_{ij \in E(\Knn)}
	\sum_{\substack{k \in \Vl, \ell \in \Vs\\k\neq i,\ell \neq j\\i\ell, k\ell, kj\in E_{g}}}(\eta_{i\ell}+\eta_{k\ell} + \eta_{kj})
\ge\frac{c_{1}}{16}\nl \ns p^3\cdot\sum_{ij\in E(\Knn)}\eta_{ij}.
\]

For fixed $k\ell\in E_{g}$, the left-hand-side may sum $\eta_{k\ell}$
as many times as the number of paths of length $3$ in $G$ that contain $k\ell$.  Each path of length $3$ can be thought of as an edge originating from $\Vl$, an edge in the middle, and an edge terminating in $\Vs$.  The total number of paths of length $3$ in $G$ containing $k\ell$ equals the number which have $k\ell$ as the middle edge, plus the number with $k\ell$ as the edge originating from $\Vl$, plus the number with $k\ell$ as the edge terminating in $\Vs$.   In each of these cases, Condition 1 ensures that there are at most $ 4 p^2 \nl \ns$ such paths of length $3$.  Hence, the term $\eta_{k\ell}$ appears at most $12 p^2 \nl \ns$ times.
Hence, the left-hand-side of above is at most 
\[
	\sum_{ij\in E(\Knn)}
	\sum_{\substack{k \in \Vl, \ell \in \Vs\\k\neq i,\ell \neq j\\i\ell, k\ell, kj\in E_{g}}}(\eta_{i\ell}+\eta_{k\ell} + \eta_{kj})
	\le 12 p^2 \nl \ns\cdot\sum_{ij\in E_{g}}\eta_{ij}.
\]
Therefore by combining the two inequalities above, we obtain 
\[
	\sum_{ij\in E(\Knn)}\eta_{ij}
	\le\frac{12\cdot 16}{c_{1}p}\sum_{ij\in E_{g}}\eta_{ij}. %\leq \sum_{ij\in E_{g}}\eta_{ij}.
	\qedhere
\]

\end{proof}

\subsection{Proof of Theorem \ref{thm:mainthm}} \label{sec-proof-mainthm}


We now prove the deterministic recovery theorem.

\begin{proof}[Proof of Theorem \ref{thm:mainthm}]

By Lemma \ref{lem:T0-feasible-bipartite} and the fact that Conditions 1--6 are invariant under global translation and nonzero scaling, we can take $\bar{\zeta}^{(0)} =0$ and $L(\Tnot, \Pnot) = 1$ without loss of generality.  The variable $\muinf =  \max_{i\neq j} \|\toij\|_2$ is to be understood accordingly.  


We will directly prove that $R(T, P) > R(\Tnot, \Pnot)$ for all $(T, P) \neq (\Tnot, \Pnot)$ such that $L(T, P)=1$ and $\Tbar + \Pbar = 0$. Consider an arbitrary feasible $T,P$ and recall the notation $\tij = (1 + \deltaij) \toij + \etaij \sij$ where $\sij$ is a unit vector orthogonal to $\toij$ and $\etaij = \| P_{\toijperp} \tij \|_2$. Since $v_{ij} =\toijhat$ holds for all $ij \in E_g$, a useful lower bound for the objective $R(T,P)$ is given by 
\begin{eqnarray}
	R(T,P)
	=\sum_{ij \in E(G)}\|\proj_{v_{ij}^{\perp}}t_{ij}\| _2
	& = & \sum_{ij\in E_{g}}\eta_{ij}+\sum_{ij\in E_{b}}\|\proj_{v_{ij}^{\perp}}t_{ij}\|_2\nonumber \\
	& \ge & \sum_{ij\in E_{g}}\eta_{ij}+\sum_{ij\in E_{b}}\left(\|\proj_{v_{ij}^{\perp}}t_{ij}^{(0)}\|_2-|\delta_{ij}|\|t_{ij}^{(0)}\|_2-\eta_{ij}\right)\nonumber \\
	& \ge & R(\Tnot, \Pnot)+\sum_{ij\in E_{g}}\eta_{ij}-\sum_{ij\in E_{b}}(|\delta_{ij}|\|t_{ij}^{(0)}\|_2+\eta_{ij}).\label{eq:gain3-bipartite}
\end{eqnarray}


Suppose that $\sum_{ij\in E_{b}}|\delta_{ij}|\|t_{ij}^{(0)}\|_2<\sum_{ij\in E_{b}}\eta_{ij}$.
Since Lemma \ref{lem:badgood-bipartite} for $\varepsilon\le\frac{c_{1}p^3}{96}$
implies $\sum_{ij\in E_{b}}\eta_{ij}\le\frac{1}{2}\sum_{ij\in E_{g}}\eta_{ij}$,
by (\ref{eq:gain3-bipartite}), we have
\begin{eqnarray*}
	R(T,P) & \ge & R(\Tnot, \Pnot)+\sum_{ij\in E_{g}}\eta_{ij}-\sum_{ij\in E_{b}}(|\delta_{ij}|\|t_{ij}^{(0)}\|_2+\eta_{ij})\\
	& > & R(\Tnot, \Pnot)+\sum_{ij\in E_{g}}\eta_{ij}-\sum_{ij\in E_{b}}2\eta_{ij}\ge R(\Tnot, \Pnot).
\end{eqnarray*}
Hence we may assume 
\begin{equation}
	\sum_{ij\in E_{b}}|\delta_{ij}|\|t_{ij}^{(0)}\|_2\ge\sum_{ij\in E_{b}}\eta_{ij}.\label{eq:badtwist-bipartite}
\end{equation}

In the case $|\Eb| \neq 0$, define $\overline{\delta}=\frac{1}{|E_{b}|}\sum_{ij\in E_{b}}|\delta_{ij}|$
as the average `relative parallel motion' on the bad edges. For a pair of vertex-disjoint edges $ij,k\ell\in E(\Knn)$, define $\eta(ij,k\ell)=\eta_{ij} + \eta_{kj} + \eta_{k\ell} + \eta_{i\ell}$,
\\

\noindent \textbf{Case 0}. $|\Eb| = 0$ or $\bar{\delta} = 0$.
\medskip

Note that $\bar{\delta}=0$ implies $\delta_{ij} =0$ for all $ij \in E_b$, which by \eqref{eq:badtwist-bipartite} implies
$\eta_{ij} =0$ for all $ij \in E_b$.  Therefore by \eqref{eq:gain3-bipartite}, we have
\[
	R(T, P) \geq R(\Tnot, \Pnot) + \sum_{ij \in E_g} \eta_{ij}.
\]
If $\sum_{ij \in E_g} \eta_{ij} > 0$, then we have $R(T, P) > R(\Tnot, \Pnot)$. Thus 
we may assume that $\eta_{ij} = 0$ for all $ij \in E_g$. In this case, we will show that $T = \Tnot$ and $P = \Pnot$.

%Therefore $\eta_{ij} = 0$ for all $ij \in E(G)$, and by
By Lemma~\ref{lem:transference-bipartite}, if $\eps \le \frac{c_1p^3}{48}$, then
$\eta_{ij}=0$ for all $ij \in E(G)$ implies that
$\eta_{ij} = 0$ for all $ij \in E(\Knn)$. 
For $ij \in E_b$, since $\delta_{ij} = \eta_{ij} = 0$, it follows that $\ell_{ij} = \ell_{ij}^{(0)}$.
Since $\delta_{ij} \|\toij\|_2 = \ell_{ij} - \ell_{ij}^{(0)}$ for $ij \in E_g$, we have
\[
	0
	= \sum_{ij \in E(G)} (\ell_{ij} - \ell_{ij}^{(0)})
	= \sum_{ij \in E_b} (\ell_{ij} - \ell_{ij}^{(0)}) + \sum_{ij \in E_g} (\ell_{ij} - \ell_{ij}^{(0)})
	= \sum_{ij \in E_g} (\ell_{ij} - \ell_{ij}^{(0)})
	= \sum_{ij \in E_g} \delta_{ij} \|\toij\|_2,
\] 
where the first equality is because $L(T, P) = L(\Tnot, \Pnot) = 1$.
By Condition 2, $\|\toij\|_2 \neq 0$ for all $i \neq j$. Therefore, if $\delta_{ij} \neq 0$ for some $ij \in E_g$, then there exists $ab, cd \in E_g$ 
such that $\delta_{ab} > 0$ and $\delta_{cd} < 0$. 
If $ab$ and $cd$ are vertex-disjoint, Lemma \ref{lem:rigidity4} and Conditions 2 and 4 force $\eta(ab,cd) > 0$, which contradicts the fact that
$\eta_{ij} = 0$ for all $ij \in E(\Knn)$.  If $ab$ and $cd$ are not vertex-disjoint, then, let $abc'd'$ be an arbitrary $C_4$ containing $ab$ and $cd$.  Then Lemma \ref{lem:rigidity4} implies the same result as above.
Therefore $\delta_{ij} = 0$ for all $ij \in E_g$, and hence
$\delta_{ij} = 0$ for all $ij \in E(G)$.  

Define $t_i = t_i^{(0)} + h_i$ for each $i \in \Vl$.  Define $p_j = \poj + h_j$ for $j \in \Vs$. Because $\eta_{ij} = \delta_{ij} = 0$ for all $ij \in E(G)$, we have $h_i = h_j$ 
for all $ij \in E(G)$. Since $G$ is connected by Condition 1, this implies $h_i = h_j$ for all 
$i \in \Vl, j \in \Vs$. Then by the constraint $\sum_{i \in \Vl} t_i + \sum_{j \in \Vs} p_j = 0$, 
we get $h_i = 0$ for all $i \in \Vl$ and $h_j = 0$ for all $j \in \Vs$. Therefore $T = \Tnot$ and $P = \Pnot$. \\

\noindent\textbf{Case 1}. $|\Eb| \neq 0$ and $\bar{\delta} \neq 0$ and $\sum_{ij\in E_{g}}|\delta_{ij}|<\frac{1}{8}\overline{\delta}|E_{g}|$. 
\medskip

Define $L_{b}=\{ij\in E_{b}:|\delta_{ij}|\ge\frac{1}{2}\overline{\delta}\}$.
Note that $\sum_{ij\in E_{b}\setminus L_{b}}|\delta_{ij}|<\frac{1}{2}\overline{\delta}|E_{b}|$
and therefore 
\begin{equation}
\sum_{ij\in L_{b}}|\delta_{ij}|=\sum_{ij\in E_{b}}|\delta_{ij}|-\sum_{ij\in E_{b}\setminus L_{b}}|\delta_{ij}|>\sum_{ij\in E_{b}}|\delta_{ij}|-\frac{1}{2}\overline{\delta}|E_{b}|=\frac{1}{2}\overline{\delta}|E_{b}|.\label{eq:large_delta-bipartite}
\end{equation}
Define $F_{g}=\{ij\in E_{g}:|\delta_{ij}|<\frac{1}{4}\overline{\delta}\}$.
Then by the condition of Case 1,
\[
\frac{1}{8}\overline{\delta}|E_{g}|>\sum_{ij\in E_{g}}|\delta_{ij}|\ge\sum_{ij\in E_{g}\setminus F_{g}}|\delta_{ij}|\ge\frac{1}{4}\overline{\delta}|E_{g}\setminus F_{g}|,
\]
and therefore $|E_{g}\setminus F_{g}|<\frac{1}{2}|E_{g}|$, or equivalently,
$|F_{g}|>\frac{1}{2}|E_{g}|$.

For each $ij\in L_{b}$, define $F_g(i,j) = \{k\ell \in F_g \mid k \neq i, \ell \neq j\}$. Note that by Condition 1, $|F_g(i,j)| > \frac{1}{2} |\Eg| - 2 p (\nl + \ns)$. For any $k\ell \in F_g(i,j)$, since $|\delta_{ij}| \ge \frac{1}{2} \overline{\delta}$ and $|\delta_{k\ell}| < \frac{1}{4}\overline{\delta}$, we have $|\delta_{ij} - \delta_{k\ell}| \ge \frac{1}{2}|\delta_{ij}|$. Thus
Lemma \ref{lem:rigidity4}
and Conditions 3 and 4 give $\eta(ij,k\ell)\ge \beta|\delta_{ij} - \delta_{k\ell}| \cdot\|\toij\|_2 \ge
\beta \cdot \frac{1}{2}|\delta_{ij}| \cdot \|\toij\|_2 \ge
\frac{\beta c_{0}\muinf}{2}|\delta_{ij}|$. 
Therefore by Condition 1, 
\begin{eqnarray*}
	\sum_{ij\in E_{b}}\sum_{\substack{k\ell\in E_{g}\\k\neq i, l \neq j}}\eta(ij,k\ell) 
	& \ge & \sum_{ij\in L_{b}}\sum_{k\ell\in F_{g}(i,j)}\frac{\beta c_{0}\muinf}{2}|\delta_{ij}|=\sum_{ij\in L_{b}}|F_{g}(i,j)|\cdot\frac{\beta c_{0}\muinf}{2}|\delta_{ij}|\\
	& >  & \sum_{ij\in L_{b}}\frac{\beta c_{0}\muinf}{2}\Bigl(\frac{1}{2}|E_{g}|  - 2p (\nl + \ns) \Bigr)|\delta_{ij}|
\end{eqnarray*}
Note that if $\eps < \frac{1}{4}p$, then $|\Eg| \geq \frac{\nl \ns p}{2} - |\Eb| \geq \frac{\nl \ns p}{4}.$
Further note that $\nl, \ns > 64$ implies that $2 p (\nl + \ns) < \frac{1}{16}\nl \ns p$.  Hence by \eqref{eq:large_delta-bipartite},
\begin{align*}
	\sum_{ij\in E_{b}}\sum_{\substack{k\ell\in E_{g}\\k\neq i, \ell \neq j}}\eta(ij,k\ell)
	& \ge\frac{\beta c_{0}\muinf}{32} \nl \ns \cdot \sum_{ij\in L_{b}} |\overline{\delta}_{ij}|
	\ge\frac{\beta c_{0}\muinf}{32} \nl \ns \cdot\frac{1}{2}\overline{\delta}|E_{b}|.
\end{align*}
For
each $ij\in E(\Knn)$, we would like to count how many times each
$\eta_{ij}$ appear on the left hand side. If $ij\in E_{b}$, then
there are at most $\nl \ns$ $C_{4}$s  containing
$ij$; hence $\eta_{ij}$ may appear at most $4 \nl \ns$
times. If $ij\notin E_{b}$, then $\eta_{ij}$ appears when there is
a $C_{4}$ containing $ij$ and some bad edge.  If the
bad edge is incident to $ij$, then there are at most
$2\eps \nl \ns$ such $C_{4}$s, and if the bad edge is
not incident to $ij$, then there are at most $|E_{b}|\le\eps \nl \ns$
such $C_{4}$s. Thus $\eta_{ij}$ may appear at most $4 \cdot 3 \eps \nl \ns = 12 \eps \nl \ns$ times. Therefore
\begin{eqnarray*}
	\sum_{ij\in E_{b}}\sum_{k\ell\in E_{g}}\eta(ij,k\ell) 
	& \le & \sum_{ij\in E_{b}}4 \nl \ns\cdot\eta_{ij}+\sum_{ij\in E(\Knn)}12 \varepsilon \nl \ns \cdot\eta_{ij}.
\end{eqnarray*}
By Lemma \ref{lem:badgood-bipartite}, if $\varepsilon < \frac{c_1p^3}{48}$, we have
\[
	\sum_{ij\in E_{b}}\sum_{k\ell\in E_{g}}\eta(ij,k\ell)
	\le\frac{48 \cdot 4\varepsilon}{c_{1}p^{3}}\nl \ns\sum_{ij\in E_{g}}\eta_{ij}+\sum_{ij\in E(\Knn)} 12 \eps \nl \ns\cdot\eta_{ij}
	\le\frac{204 \varepsilon}{c_{1}p^{3}}\nl \ns\sum_{ij\in E(\Knn)}\eta_{ij}.
\]
Hence
\[
	\frac{204\varepsilon}{c_{1}p^{3}}\nl \ns\sum_{ij\in E(\Knn)}\eta_{ij}\ge\frac{\beta c_{0}\muinf}{64} \nl \ns \cdot\overline{\delta}|E_{b}|.
\]

If $\varepsilon< \frac{\beta c_0 c_1^2 p^4}{384\cdot 204 \cdot 64}$, then by Condition 3,
 $\bar{\delta} \neq 0$, and $|\Eb| \neq 0$, the above implies 
\begin{align*}
	\sum_{ij\in E(\Knn)}\eta_{ij}
	\ge&\, \frac{\beta c_{0}c_{1}p^{3}}{204 \cdot 64 \eps}\muinf \cdot\overline{\delta}|E_{b}| \\
	%\ge \frac{\beta c_{0}c_{1}p^{3}}{42\cdot 32 \cdot 8} \cdot \frac{1}{\varepsilon} \cdot \muinf \overline{\delta}|E_{b}|  \\
	>&\, \frac{384}{c_{1}p}\muinf\cdot\overline{\delta}|E_{b}|
	\ge \frac{384}{c_{1}p}\sum_{ij\in E_{b}}|\delta_{ij}|\|t_{ij}^{(0)}\|_2.
\end{align*}
Lemma \ref{lem:transference-bipartite} implies 
\[
	\sum_{ij\in E_{g}}\eta_{ij}
	\ge \frac{c_{1}p}{192}\sum_{ij\in E(\Knn)}\eta_{ij}
	> 2\sum_{ij\in E_{b}}|\delta_{ij}|\|t_{ij}^{(0)}\|_2.
\]
%Lemma \ref{lem:badgood} implies $\sum_{ij\in E_{g}}\eta_{ij}\ge 2\sum_{ij\in E_{b}}\eta_{ij}$ if $\varepsilon < \frac{c_1p^2}{16}$.
Therefore by (\ref{eq:badtwist-bipartite}),we have $\sum_{ij\in E_{g}}\eta_{ij}>\sum_{ij\in E_{b}}(|\delta_{ij}|\|t_{ij}^{(0)}\|_2+\eta_{ij})$ if $\eps \le \min\{\frac{c_1p^3}{48}, \frac{p}{4}, \frac{\beta c_0 c_1^2 p^4}{384 \cdot 204 \cdot 64} \}$ and $p \geq \sqrt{\frac{8}{n}}$.
By (\ref{eq:gain3-bipartite}), this shows $R(T,P)>R(\Tnot, \Pnot)$. This condition on $\eps$ is satisfied under the assumption $\eps \leq \frac{\beta c_0 c_1^2 p^4}{384 \cdot 204\cdot64}$.
\\

\noindent\textbf{Case 2}. $|\Eb| \neq 0$ and $\bar{\delta} \neq 0$ and $\sum_{ij\in E_{g}}|\delta_{ij}|\ge\frac{1}{8}\overline{\delta}|E_{g}|$. 
\medskip

Define $E_{+}=\{ij\in E_{g}\,:\,\delta_{ij}\ge0\}$ and $E_{-}=\{ij\in E_{g}\,:\,\delta_{ij}<0\}$.
Since $\ell_{ij}-\ell_{ij}^{(0)}=\delta_{ij}\|t_{ij}^{(0)}\|_2$
for $ij\in E_{g}$, we have 
\begin{eqnarray*}
0=\sum_{ij \in E(G)}(\ell_{ij}-\ell_{ij}^{(0)}) & = & \sum_{ij\in E_{b}}(\ell_{ij}-\ell_{ij}^{(0)})+\sum_{ij\in E_{g}}\delta_{ij}\|t_{ij}^{(0)}\|_2.
\end{eqnarray*}
where the first equality follows from $L(T,P) = L(\Tnot, \Pnot)$.  Therefore, 
\begin{eqnarray*}
	\left| \sum_{ij\in E_{g}}\delta_{ij}\|t_{ij}^{(0)}\|_2 \right|
	\le \left| \sum_{ij\in E_{b}}(\ell_{ij}-\ell_{ij}^{(0)}) \right|
	\le \sum_{ij\in E_{b}}(|\delta_{ij}|\|t_{ij}^{(0)}\|_2+\eta_{ij}) 
	\le 2\muinf\overline{\delta}|E_{b}|,
\end{eqnarray*}
where the last inequality follows from (\ref{eq:badtwist-bipartite}), Condition 3, and the definition of $\overline{\delta}$. On the other hand, the condition of Case 2 and Condition 3 implies 
$\sum_{ij\in E_{g}}|\delta_{ij}|\|t_{ij}^{(0)}\|_2 \ge \frac{1}{8}c_0 \muinf\overline{\delta}|E_g|$. 
Therefore
\[
	\sum_{ij \in E_-} (-\delta_{ij}) \|t_{ij}^{(0)}\|_2
	= \frac{1}{2}\left( -\sum_{ij\in E_{g}}\delta_{ij}\|t_{ij}^{(0)}\|_2 + \sum_{ij\in E_{g}}|\delta_{ij}|\|t_{ij}^{(0)}\|_2\right)
	\ge \frac{1}{2} \left(\frac{1}{8} c_0 \muinf\overline{\delta}|E_g| - 2\muinf\overline{\delta}|E_{b}|\right).
\]
If $\varepsilon \le \frac{1}{128}c_0p$, then 
since $|E_b| \le \varepsilon \nl \ns$ and $|E_g| \ge \frac{\nl \ns p}{2} - |E_b| \ge \frac{\nl \ns p}{4}$, 
we see that $\frac{1}{8} c_0 \muinf\overline{\delta}|E_g| - 2\muinf\overline{\delta}|E_{b}| \ge \frac{1}{16}c_0\muinf \bar{\delta}|E_g|$.
Therefore $\sum_{ij\in E_{-}}(-\delta_{ij})\|t_{ij}^{(0)}\|_2 \ge\frac{1}{32}c_{0}\muinf \overline{\delta}|E_{g}|$.
Similarly, $\sum_{ij\in E_{+}} \delta_{ij}\|t_{ij}^{(0)}\|_2 \ge \frac{1}{32}c_{0}\muinf \overline{\delta}|E_{g}|$.

If $|E_+| \ge \frac{1}{2}|E_g|$, then 
by Lemma \ref{lem:rigidity4} and Conditions 3 and 4, we have 
\begin{eqnarray*}
	\sum_{ij\in E_{-}}\sum_{\substack{k\ell\in E_{+}\\k\neq i, \ell \neq j}}\eta(ij,k\ell) 
	& \ge & \sum_{ij\in E_{-}}\sum_{\substack{k\ell\in E_{+}\\k\neq i, \ell \neq j}}\beta(-\delta_{ij})\|t_{ij}^{(0)}\|_2\\
	& \ge & \sum_{ij\in E_{-}}(-\delta_{ij})\|t_{ij}^{(0)}\|_2\cdot \beta(|E_{+}| - 2 p (\nl + \ns) ) \\
	& \ge & \frac{1}{32}c_{0}\muinf \overline{\delta}|E_{g}| \cdot \beta (|E_{+}| - 2 p (\nl + \ns) ) \\
	& \ge & \frac{\beta}{32} c_{0} \muinf\overline{\delta}|E_{g}| \Bigl(\frac{1}{2} | E_g| - 2 p (\nl + \ns) \Bigr).
\end{eqnarray*}
Note that if $\eps < \frac{1}{4}p$, then $|\Eg| \geq \frac{\nl \ns p}{2} - |\Eb| \geq \frac{\nl \ns p}{4}.$
Further note that $\nl, \ns > 64$ implies that $2 p (\nl + \ns) < \frac{1}{16}\nl \ns p$.  Hence,
\begin{eqnarray*}
	\sum_{ij\in E_{-}}\sum_{\substack{k\ell\in E_{+}\\k\neq i, l \neq j}}\eta(ij,k\ell) 
	&\ge & 
	\frac{1}{32}\beta c_0 \muinf \overline{\delta}\cdot \frac{\nl \ns p}{4} \cdot \frac{\nl \ns p}{16}
	\ge
	\frac{\beta c_0 \muinf \overline{\delta} \nl^2 \ns^2 p^2}{32 \cdot 64}.
\end{eqnarray*}
Similarly, if $|E_-| \ge \frac{1}{2}|E_g|$, then we can switch the order of summation and
consider $\sum_{ij \in E_+} \sum_{k\ell \in E_-} \eta(ij, k\ell)$ to obtain the 
same conclusion.

Since each edge is contained in at most $\nl \ns$
copies of $C_{4}$ and there are 4 edges
in a $C_{4}$, we have 
\[
	\sum_{ij\in E_{-}}\sum_{\substack{k\ell\in E_{+}\\k\neq i, \ell \neq j}}\eta(ij,k\ell)
	\le 4 \nl \ns \sum_{ij\in E(\Knn)}\eta_{ij}.
\]
If $\varepsilon < \frac{\beta c_0 c_1 p^3}{384\cdot4\cdot32\cdot64}$, then since $\bar{\delta} \neq 0$ and $|\Eb| \leq \eps \nl \ns$, we have
\[
	\sum_{ij\in E(\Knn)}\eta_{ij}
	\ge\frac{1}{4 \nl \ns}\cdot\frac{\beta c_{0} \muinf \overline{\delta}}{32\cdot 64} \nl^2 \ns^2 p^2
	\ge\frac{\beta c_{0} p^2}{4 \cdot 32 \cdot 64}\muinf \overline{\delta} \nl \ns
	>\frac{384}{c_{1}p}\muinf\overline{\delta}|E_{b}|.
\]
By Lemma \ref{lem:transference-bipartite}, if $\varepsilon < \frac{c_1p^3}{48}$, then this implies 
\[
\sum_{ij\in E_{g}}\eta_{ij}\ge\frac{c_{1}p}{192}\sum_{ij\in E(\Knn)}\eta_{ij} > 2\muinf\overline{\delta}|E_{b}|.
\]
Therefore from (\ref{eq:gain3-bipartite}), (\ref{eq:badtwist-bipartite}), and Condition 3, 
if $\varepsilon \le \min\{\frac{c_0 p}{128}, \frac{c_1p^3}{96}, \frac{p}{4}, \frac{\beta c_0 c_1 p^3}{384\cdot4\cdot32\cdot64} \}$ and $p \geq \sqrt{\frac{8}{n}}$, then
\begin{eqnarray*}
R(T, P) & \ge & R(\Tnot, \Pnot)+\sum_{ij\in E_{g}}\eta_{ij}-\sum_{ij\in E_{b}}(|\delta_{ij}|\|t_{ij}^{(0)}\|_2+\eta_{ij})\\
 & > & R(\Tnot, \Pnot)+2\muinf\overline{\delta}|E_{b}|-\sum_{ij\in E_{b}}2|\delta_{ij}|\|t_{ij}^{(0)}\|_2\ge R(\Tnot, \Pnot).
\end{eqnarray*}
This condition on $\eps$ is satisfied under the assumption $\eps \leq \frac{\beta c_0 c_1^2 p^4}{384\cdot204\cdot64}$.
\end{proof}


\subsection{Properties of Gaussians} \label{sec-properties-of-gaussians}

In this section, we prove that i.i.d. Gaussians satisfy properties needed to establish Conditions 3--5 in Theorem \ref{thm:mainthm}. We begin by recording some useful facts regarding concentration of random Gaussian vectors:

\begin{lemma}
\label{lem:conc1}
Let $x,y$ be i.i.d. $\mathcal{N}(0,I_{d\times d})$, and $\epsilon \leq 1$, then
\[
\P\left ( d(1-\epsilon) \leq \|x\|_2^2 \leq d(1+\epsilon) \right) \geq 1- e^{-c \epsilon^2 d}
\]
and
\[
\P \left( |\langle x, y \rangle| \geq d\epsilon \right) \leq e^{-c\epsilon^2 d}
\]
where $c>0$ is an absolute constant.
\end{lemma}
\begin{proof}
Both statements follow from Corollary 5.17 in \cite{Vershynin}, concerning concentration of sub-exponential random variables.
\end{proof}

\begin{lemma}
\label{lem:conc2}
Corollary 5.35 in \cite{Vershynin}. Let $A$ be an $n \times d$ matrix with iid $\mathcal{N}(0,1)$ entries. Then for any $t\geq 0$,
\[
\P\left( \sigma_{\text{max}}(A) \geq \sqrt{n} + \sqrt{d} + t \right) \leq 2e^{-\frac{t^2}{2}}
\]
where $\sigma_{\max}(A)$ is the largest singular value of $A$.
\end{lemma}


\begin{lemma} \label{lem-beta-first}
Let $\ti, \tj, \tk, \tl \sim \mathcal{N}(0, I_{d \times d})$ be independent.  There is a universal constant $c$ such that with probability at least $1 - 15 e^{-cd}$,
$$
\frac{\|\proj_{\Span(\tk - \tj, \ti - \tl)^\perp} (\ti - \tj) \|}{\| \ti - \tj\|} \geq \frac{1}{4}.
$$
\end{lemma}
\begin{proof}
Let $c$ be the constant from Lemma \ref{lem:conc1}.
Let $x = \ti - \tl, y = \tk -\tj , z = \ti - \tj$.  Observe
\begin{align*}
\proj_{\Span(x,y)^\perp} \zhat  
= \zhat - \langle \zhat, \xhat \rangle \xhat - \langle \zhat, \yxphat \rangle \yxphat 
=&\, \zhat - \langle \zhat, \xhat \rangle \xhat - \langle \zhat, \yhat \rangle \yhat + (\langle \zhat, \yhat \rangle \yhat - \langle \zhat, \yxphat \rangle \yxphat) \\
  =&\, \zhat - \langle \zhat, \xhat \rangle \xhat - \langle \zhat, \yhat \rangle \yhat +(\yhat \yhat^t - \yxphat \yxphat^t ) \zhat,
\end{align*}
where $\yxphat = \frac{y - \langle y, \xhat \rangle \xhat}{\|y - \langle y, \xhat \rangle \xhat \|}$, which is well defined with probability $1$.  By the triangle inequality,
$$
\| \proj_{\Span(x,y)^\perp} \zhat \| \geq \sqrt{1 - | \langle \zhat, \xhat \rangle |^2} - | \langle \zhat, \yhat \rangle |  - \| \yhat \yhat^t - \yxphat \yxphat^t \|_{\text{op}}
$$
For arbitrary unit vectors $\ahat, \bhat \in \R^d$,  $\| \ahat \ahat^t - \bhat \bhat^t\|_\text{op} =  |\sin \theta|$, where $\theta$ is the angle between $\ahat$ and $\bhat$.  This fact can be verified by direct computation after taking $\ahat = e_1$ and $\bhat = \cos \theta \ e_1 + \sin \theta \ e_2$ without loss of generality.  Hence, $\| \yhat \yhat^t - \yxphat \yxphat^t\|_\text{op} =  | \sin \theta | =  | \cos \alpha |$, where $\theta$ is the angle between $\yhat$ and $\yxphat$, and $\alpha$ is the angle between $\yhat$ and $\xhat$.  Thus $\| \yhat \yhat^t - \yxphat \yxphat^t\|_\text{op} = | \langle \yhat, \xhat \rangle|$. 
So, 
$$
\| \proj_{\Span(x,y)^\perp} \zhat \| \geq \sqrt{1 - | \langle \zhat, \xhat \rangle |^2} - | \langle \zhat, \yhat \rangle | - | \langle \yhat, \xhat \rangle |
$$
Now, note that 
$$
| \langle \zhat, \xhat \rangle |^2 = \frac{\langle \ti - \tj, \ti - \tl \rangle^2}{\| \ti - \tj \|^2 \| \ti - \tl \|^2} = \frac{(\|\ti\|^2 - \langle \ti, \tl \rangle - \langle \tj, \ti \rangle + \langle \tj, \tl \rangle)^2}{\| \ti - \tj \|^2 \| \ti - \tl \|^2}
$$
By Lemma \ref{lem:conc1} with $\eps = 0.01$,
$$
| \langle \zhat, \xhat  \rangle |^2 \leq \frac{(d ( 1+ \eps) + 3d \eps )^2}{4 d^2 (1-\eps)^2} \leq 0.3
$$
with probability at least $1 - 6 e^{-c d}$ for some universal constant $c$.   Similarly, $| \langle \zhat, \yhat \rangle |^2 \leq 0.3$ with the same probability.
Since $\yhat$ and $\xhat$ are independent, by Lemma \ref{lem:conc1} with $\eps = 0.01$, $| \langle \yhat, \xhat \rangle | \leq \frac{\eps d}{d (1-\eps)} \leq 2 \eps$ with probability at least $1 - 3 e^{-c d}$.
Thus, we observe 
$$
\| \proj_{\Span(x,y)^\perp} \zhat \| \geq \sqrt{1 - 0.3} - \sqrt{0.3} - 0.02 \geq \frac{1}{4}
$$
with probability at least $1 - 15 e^{-c d}$. 
\end{proof}

\begin{lemma} \label{lem-beta-second}
Let $\ti, \tj, \tk, \tl \sim \mathcal{N}(0, I_{d \times d})$ be independent for $d \geq 3$.  There is a universal constant $c$ such that with probability at least $1 - 7 e^{-cd}$,
$$
\frac{\|\proj_{\Span(\ti - \tl, \tk - \tl)^\perp} (\ti - \tj) \|}{\|\ti - \tj \|_2} \geq \frac{1}{4}.
$$
\end{lemma}
\begin{proof}
%Note that $\| \ti - \tj\|^2 \leq 2d (1+\eps)$ with probability at least $1 - e^{-c d}$ for $\eps = 0.01$.  Thus it suffices to show that 
%$\|\proj_{\Span(\tj - \tk, \tj - \tl)^\perp} (\ti - \tj) \| $.
Let $c$ be the constant from Lemma \ref{lem:conc1}.
Let $u = \frac{\ti - \tl}{\sqrt{2}}, v = \frac{\tl + \ti}{\sqrt{2}}, w = \tj, x = \tk$.  Each of these variables are i.i.d. $\mathcal{N}(0, 1)$.
Note that 
$$
\proj_{\Span(\ti - \tl, \tk - \tl)^\perp} (\ti - \tj) = -\proj_{\Span(u, x - \frac{v}{\sqrt{2}})^\perp} \Bigl(w - \frac{v}{\sqrt{2}} \Bigr).
$$
Without loss of generality, rotate coordinates so that $u$ is in the direction of $e_1$.  Thus, it suffices to bound
$
\bigl \| \proj_{(\xtilde - \frac{\vtilde}{\sqrt{2}})^\perp} ( \wtilde - \frac{\vtilde}{\sqrt{2}} ) \bigr \|_2
$
where $\vtilde, \wtilde, \xtilde \sim \mathcal{N}(0, I_{d-1 \times d-1})$.  Note that $\xtilde - \frac{\vtilde}{\sqrt{2}}$ and $\wtilde - \frac{\vtilde}{\sqrt{2}}$ both follow the distribution $\mathcal{N}(0,\frac{3}{2} I_{d-1 \times d-1})$.
Note that 
\begin{align*}
\Bigl \| \proj_{(\xtilde - \frac{\vtilde}{\sqrt{2}})^\perp} \Bigl( \wtilde - \frac{\vtilde}{\sqrt{2}} \Bigr) \Bigr \|^2 
&= \Bigl \|\wtilde - \frac{\vtilde}{\sqrt{2}} \Bigr \|^2 - \frac{\langle \wtilde - \frac{\vtilde}{\sqrt{2}}, \xtilde - \frac{\vtilde}{\sqrt{2}} \rangle^2}{\| \xtilde - \frac{\vtilde}{\sqrt{2}}\|^2} \\
&= \Bigl \|\wtilde - \frac{\vtilde}{\sqrt{2}} \Bigr \|^2 - \frac{\bigl (\langle \wtilde, \xtilde \rangle - \frac{1}{\sqrt{2}} \langle \wtilde, \vtilde \rangle - \frac{1}{\sqrt{2}} \langle \vtilde, \xtilde \rangle + \frac{\| \vtilde\|^2}{2} \bigr)^2}{\| \xtilde - \frac{\vtilde}{\sqrt{2}}\|^2}.
\end{align*}

Hence Lemma \ref{lem:conc1} with $\eps = 0.01$ shows that  with probability at least $1 - 6 e^{-c d}$, the above is at least
\begin{align*}
\frac{3}{2} (d-1) (1-\eps) - \frac{\Bigl(\frac{1}{2} (d-1) (1+\eps) - 3 \eps (d-1)\Bigr)^2}{\frac{3}{2} (d-1) (1-\eps)} \geq (d-1) \geq \frac{2}{3} d,
\end{align*}
%where the first inequality follows from Lemma \ref{lem:conc1} with $\eps = 0.01$ and holds with probability at least $1 - 6 e^{-c d}$.
Thus, we have that 
$$
\|\proj_{\Span(\tj - \tl, \tk - \tl)^\perp} (\ti - \tj) \|^2  \geq \frac{2}{3} d
$$
with probability at least $1 - 6 e^{-c d}$.  To conclude the proof, note that Lemma \ref{lem:conc1} with $\eps = 0.01$ implies that $\|\ti - \tj\|^2 \geq 2 d ( 1+\eps)$ with probability at least $1 - e^{-c d}$.
\end{proof}

We can now establish Conditions 3--4 of Theorem~\ref{thm:mainthm} with high probability.
\begin{lemma}\label{lem-conds34}
Let $\ti, \pj \sim \mathcal{N}(0, I_{d \times d})$ for $i \in \Vl, j \in \Vs$ be independent.   Condition 3 of Theorem~\ref{thm:mainthm} holds with $c_0 = \frac{9}{10}$ with probability at least $1 - 2 \nl \ns e^{-cd}$ for a universal constant $c$.  Condition 4 of Theorem~\ref{thm:mainthm} holds with $\beta = \frac{1}{4}$ with probability at least $ 1 - 22 \nl^2 \ns^2 e^{-c d}$. 
\end{lemma}
\begin{proof}
Condition 3 follows from applying Lemma \ref{lem:conc1} with $\eps=0.01$ and a union bound to $\|\ti - \pj\|_2$ for all $\nl \ns$ pairs $(i,j) \in \Vl \times \Vs$.  Condition 4 follows from applying Lemmas \ref{lem-beta-first} and \ref{lem-beta-second} and a union bound over the at most $\nl^2 \ns^2$ choices of $i, k \in \Vl$ and $j, \ell \in \Vs$. 
\end{proof}


\subsection{Gaussians are well-distributed} \label{sec-gaussians-well-distributed}

In this section we prove that Condition 5 of Theorem~\ref{thm:mainthm} holds with high probability.

\begin{lemma} \label{lem-c4-well-distributed}
There exist constants $d_0$ and $K_0$ such that the following holds.  Let $G = (\Vl \cup \Vs, E)$ be a bipartite-$p$-typical graph.  Let $\ti, \pj \sim \mathcal{N}(0, I_{d \times d})$ for $i \in \Vl, j \in \Vs$ be independent from $G$ and each other.  Let $T = \{\ti\}_{i \in \Vl}$ and $P = \{\pj\}_{j \in \Vs}$.  If $d \geq d_0$ and $\nl, \ns \geq \max(K_0, 160 d)$, then $(T, P)$ is $\frac{1}{20}$-well-distributed along $G$ with probability at least $1 - O( \nl^2 \ns^2 e^{-cd})$ for  universal constants $c, K_0$.
\end{lemma}

We start by proving an intermediate lemma asserting the well-distributedness
of pairs of random Gaussian vectors $\{(t_i, p_i)\}_{i \in [k]}$ with respect to a fixed pair of random Gaussian vectors $(x,y)$.

\begin{lemma} \label{lem:wd_over_matching}
There exist positive constants $d_0,\Kotilde$ such that the following holds. Let $x,y, \ti, p_i \sim \mathcal{N}(0, I_{d \times d})$ be independent, where $i \in [k]$. Then the set $\{(t_i, p_i)\}_{i \in [k]}$
is $\frac{1}{10}$-well-distributed with respect to $(x,y)$ with probability $1 - 6 k e^{-cd}$ if $k \geq \max(\Kotilde, 10d)$ and $d \geq d_0$.
\end{lemma}

The proof of this lemma appears at the end of this section.
We will deduce Lemma~\ref{lem-c4-well-distributed} from Lemma~\ref{lem:wd_over_matching}
by partitioning the edge set of $G$ into sets of vertex-disjoint edges.
A {\em matching} is a set of vertex-disjoint edges.
A {\em perfect matching} of a graph is a matching that intersects all vertices.
The following is a well-known lemma in Graph theory.

\begin{lemma} \label{lem:decompose}
Let $G = (V, E)$ be a bipartite graph with vertex partition $V = V_1 \cup V_2$, and
let $\Delta$ be the maximum degree of $G$. 
There exists an edge-partition $E = E_1 \cup \cdots \cup E_\Delta$ such that $E_a$ forms a matching
for each $a \in [\Delta]$.
\end{lemma}
\begin{proof}
By adding vertices and edges to $G$ if necessary, we can obtain a $\Delta$-regular 
bipartite multi-graph $G'$.
By Hall's theorem, every non-empty
regular multi-graph contains a perfect matching (see \cite[Corollary 2.1.3]{Diestel}). 
Let $F_1$ be an arbitrary perfect matching of $G'$. Remove $F_1$ from the edge set of $G'$, 
and note that the remaining graph is still regular. Thus we can repeat the process to 
obtain a partition $E(G') = F_1 \cup \cdots \cup F_{\Delta}$ of the edge set of $G'$ into perfect 
matchings. 
The sets $E_a = F_a \cap E(G)$ for $a \in [\Delta]$ satisfy the claimed condition.
\end{proof}

The proof of Lemma~\ref{lem-c4-well-distributed} follows from the two lemmas above.

\begin{proof}[Proof of Lemma~\ref{lem-c4-well-distributed}]
Recall the notation that $N = \max\{|\Vl|, |\Vs|\}$ and $n = \min\{|\Vl|, |\Vs|\}$.  Since $G$ is a bipartite-$p$-typical graph, the maximum degree $\Delta$ of $G$ is at most
$2N p$. By Lemma~\ref{lem:decompose}, there exists an edge-partition 
$E = E_1 \cup \cdots \cup E_\Delta$ such that each $E_a$ for $a=1,2,\ldots, \Delta$ forms
a matching.

Fix a pair of indices $(i_0,j_0)$ for $i_0 \in \Vl$ and $j_0 \in \Vs$.
Let $E' \subseteq E$ be the subset of edges that do not intersect $i_{0}$ or $j_{0}$, and
for each $a \in [\Delta]$, let $E_a' \subseteq E_a$ be the subset of edges that
do not intersect $i_{0}$ or $j_{0}$.
Let $A \subseteq [\Delta]$ be the set of indices $a$ for which $|E_a'| \ge \max(\tilde{K}_0, 10 d)$.
For each $a \in A$, by Lemma~\ref{lem:wd_over_matching}, we see that 
with probability at least $1 - O(|E_a'| e^{-cd})$,
\[
	\sum_{ij \in E_a'} \| \proj_{\Span\{p_{j_0}-t_{i},t_{i}-p_{j}, p_j-t_{i_0}\}^{\perp}} (h) \|_2
	\ge \frac{1}{10} |E_a'| \| h \|_2
\]
holds for all $h \in \mathbb{R}^d$.
Therefore by the union bound, with probability at least $1 - O(\sum_{a \in A} |E_a'| e^{-cd}) \geq 1 - O( |E'| e^{-cd}) \geq 1 - O( nN e^{-cd})$, the
above holds simultaneously for all $a \in [\Delta]$. Conditioned on this event, for all $h \in \mathbb{R}^d$,
\begin{align*}
	&\, \sum_{ij \in E'} \| \proj_{\Span\{p_{j_0}-t_{i},t_{i}-p_{j}, p_j-t_{i_0}\}^{\perp}} (h) \|_2 \\
	\ge &\,
	\sum_{a \in A} \sum_{ij \in E_a'} \| \proj_{\Span\{p_{j_0}-t_{i},t_{i}-p_{j}, p_j-t_{i_0}\}^{\perp}} (h) \|_2
	\ge \sum_{a \in A} \frac{1}{10} |E_a'| \| h \|_2.
\end{align*}
Since $|E'| = \sum_{a =1}^{\Delta} |E_a'|$, we see that
\begin{align} \label{eq:combine_rotation}
	\sum_{ij \in E'} \| \proj_{\Span\{p_{j_0}-t_{i},t_{i}-p_{j}, p_j-t_{i_0}\}^{\perp}} (h) \|_2
	\ge \frac{1}{10} \left( |E'| - \sum_{a \notin A} |E_a'| \right) \| h \|_2.
\end{align}
Since $G$ is bipartite-$p$-typical, we have $|E'| \ge \frac{1}{2}nNp - 2(N+n)p \ge \frac{1}{4}nNp$ if $n > 16$, and by the  definition of $A$, we have
$\sum_{a \notin A} |E_a'| \le  \max(\tilde{K}_0, 10 d) \cdot \Delta \le \frac{1}{8}nNp$ if $n \geq 16 \cdot \max(\tilde{K}_0, 10 d)$.
Hence the right-hand-side of \eqref{eq:combine_rotation} is at least
$\frac{1}{20}|E'| \| h \|_2$ for all $h \in \mathbb{R}^d$. This shows that the set $\{(\ti,\pj)\}_{i \neq i_0, j \neq j_0}$
is $\frac{1}{20}$-well-distributed with respect to $(t_{i_0}, p_{j_0})$ with probability 
at least $1 - O(n N e^{-cd})$. By taking the union bound over all 
choices of pairs $(i_0, j_0) \in \Vl \times \Vs$, we can conclude that $(T,P)$ is
$\frac{1}{20}$-well-distributed along $G$ with probability at least 
$1 - O(n^2 N^2 e^{-cd})$.
\end{proof}

We now prove Lemma~\ref{lem:wd_over_matching}.

\begin{proof}[Proof of Lemma~\ref{lem:wd_over_matching}]

Throughout the proof, the positive constant $c$ may change from line to line, but is always bounded below by the positive constant of the lemma statement.

For each $i$, let $
W_i = \Span(t_i - y, p_i -x, t_i - p_i) =\Span( x-y, p_i + t_i - (x+y), t_i-p_i).
$
Thus $P_{W_i^\perp} \circ P_{(x-y)^\perp} = P_{W_i^\perp}$. Therefore, it is enough to show that for all $h \perp x-y$, with high probability
\[
\sum_{i=1}^n \| P_{W_i^\perp}( h) \|_2 \geq \frac{1}{10}n \|h\|_2.
\]

%Now, we have
%\[
%P_{(t_i-p_i)^\perp}( 2p_i - (x+y) ) \approx 2p_i - (x+y) - \langle 2p_i - (x+y), t_i-p_i \rangle (t_i -p_i) \frac{1}{2d} \approx 2p_i - (x+y) + t_i-p_i = p_i + t_i - (x+y)
%\]
%where we used that $\|t_i\|_2^2 \approx d$ and that two independent gaussian vectors are approximately orthogonal. Therefore,
Letting $V_i = \Span(x-y, p_i + t_i - x -y)$, we have
\[
W_i = \Span( x-y, p_i + t_i - x-y, t_i-p_i) = \Span \Bigl( x-y, p_i + t_i - x-y, P_{V_i^\perp}(t_i-p_i) \Bigr).
\]
Now, for any $h \perp (x-y)$, 
\begin{align*}
\sum_{i=1}^n \| P_{W_i^\perp}(h)\|_2  
& \geq \left\| \sum_{i=1}^n P_{W_i^\perp}( h) \right\|_2 \\
& = \left\| \sum_{i=1}^n \left(P_{V_i^\perp}( h) - P_{P_{V_i^\perp}(t_i-p_i)}(h) \right) \right\|_2 \\
& \geq \left\| \sum_{i=1}^n P_{V_i^\perp}( h)\right\|_2  - \left\| \sum_{i=1}^n P_{P_{V_i^\perp}(t_i-p_i)}(h)  \right\|_2 \\
& \geq \left\| \sum_{i=1}^n P_{V_i^\perp}( h)\right\|_2  - \sum_{i=1}^n \left\| P_{P_{V_i^\perp}(t_i-p_i)}(h) - P_{(t_i-p_i)}(h) \right\|_2 - \left\|\sum_{i=1}^n P_{(t_i-p_i)}(h)\right\|_2.
\end{align*}

Since $ \|P_v(h) - P_w (h) \|_2 \leq \| \vhat \vhat^t - \what \what^t\|_\text{op} \|h\|_2  \leq  \| \hat{v} - \hat{w} \|_2 \|h\|_2 $ holds for all vectors $v,w,h \in \mathbb{R}^d$, the above is at least
\begin{align} \label{eq:temp_midstep}
& \left\| \sum_{i=1}^n P_{V_i^\perp}( h)\right\|_2 - \|h\|_2 \sum_{i=1}^n \left \| \frac{P_{V_i^\perp}(t_i-p_i)}{\|P_{V_i^\perp}(t_i-p_i)\|_2} - \frac{(t_i-p_i)}{\|t_i - p_i\|_2}  \right \|_2 - \left\|\sum_{i=1}^n P_{(t_i-p_i)}(h)\right\|_2.
\end{align}
Note that when $v = P(w)$ for some orthogonal projection operator $P$, we have
\[
\| \hat{v} - \hat{w}\|_2^2 = 2\left(1-\langle \hat{v},\hat{w} \rangle\right) = 2\left(1 - \frac{\langle  P(w), w \rangle}{\|P(w)\|_2 \|w\|_2}\right) = 2\left(1- \frac{\|P(w)\|_2}{\|w\|_2}\right).
\]
Thus,
\[
\left \| \frac{P_{V_i^\perp}(t_i-p_i)}{\|P_{V_i^\perp}(t_i-p_i)\|_2} - \frac{(t_i-p_i)}{\|t_i - p_i\|_2}  \right \|_2 = \sqrt{2}\sqrt{1 - \frac{\| P_{V_i^\perp} (t_i-p_i) \|_2}{\|t_i-p_i\|_2}}.
\]
Hence \eqref{eq:temp_midstep} is at least
\begin{align*}
& \left\| \sum_{i=1}^n P_{V_i^\perp}( h)\right\|_2 -  \|h\|_2 \sum_{i=1}^n \left \| \frac{P_{V_i^\perp}(t_i-p_i)}{\|P_{V_i^\perp}(t_i-p_i)\|_2} - \frac{(t_i-p_i)}{\|t_i - p_i\|_2}  \right \|_2 - \frac{1}{\min_i(\|t_i - p_i\|_2^2)}\left \|\sum_{i=1}^n (t_i-p_i)(t_i-p_i)^*\right\|_{\text{op}} \|h\|_2\\
&= \left\| \sum_{i=1}^n P_{V_i^\perp}( h)\right\|_2 -  \|h\|_2 \sum_{i=1}^n \sqrt{2}\sqrt{1 - \frac{\| P_{V_i^\perp} (t_i-p_i) \|_2}{\|t_i-p_i\|_2}} - \frac{1}{\min_i(\|t_i - p_i\|_2^2)}\left \|\sum_{i=1}^n (t_i-p_i)(t_i-p_i)^*\right\|_{\text{op}} \|h\|_2.
%& \geq c n \|h\|_2 - \| \sum_{i=1}^n P_{P_{V_i^\perp}(t_i-p_i)}(h)  \|_2 
\end{align*}
%where the \red{third} inequality follows from the fact that $ \|P_v(h) - P_w (h) \|_2 \leq \| \vhat \vhat^t - \what \what^t\|_\text{op} \|h\|_2  \leq \red{1} \| \hat{v} - \hat{w} \|_2 \|h\|_2 $, and 
%the last equality follows from the fact that when $v = P(w)$ for some orthogonal projection operator $P$, we have
%\[
%\| \hat{v} - \hat{w}\|_2^2 = 2\left(1-\langle \hat{v},\hat{w} \rangle\right) = 2\left(1 - \frac{\langle  P(w), w \rangle}{\|P(w)\|_2 \|w\|_2}\right) = 2\left(1- \frac{\|P(w)\|_2}{\|w\|_2}\right).
%\]
We will now expand the first term, $\| \sum_{i=1}^n P_{V_i^\perp}( h)\|_2$. Let $u = x-y$ and $z_i = x+y - (p_i + t_i)$. We have
\begin{align*}
P_{V_i^\perp}( h) &= P_{S(u,z_i)^\perp}(h)\\
&= P_{S(u,P_{u^\perp} (z_i))^\perp} (h) \\
& = h - P_{u}(h) - P_{P_{u^\perp} (z_i)}(h) \\
& = h - P_{P_{u^\perp} (z_i)}(h) + P_{z_i} (h) - P_{z_i} (h)\\
& =  P_{z_i^\perp} (h) - P_{P_{u^\perp} (z_i)}(h) + P_{z_i} (h),
\end{align*}
where we used $h \perp u$ in the fourth inequality. 
Thus,
\begin{align*}
\| \sum_{i=1}^n P_{V_i^\perp}( h)\|_2 &= \| \sum_{i=1}^n \left(P_{z_i^\perp}(h)  - P_{P_{u^\perp} (z_i)}(h) + P_{z_i} (h) \right) \|_2 \\
&\geq \| \sum_{i=1}^n P_{z_i^\perp}(h)\|_2 - \sum_{i=1}^n \| P_{P_{u^\perp} (z_i)}(h) - P_{z_i} (h) \|_2 \\
& \geq \| \sum_{i=1}^n P_{z_i^\perp}(h)\|_2 - \|h\|_2 \sum_{i=1}^n \left \| \frac{P_{u^\perp} (z_i)}{\|P_{u^\perp} (z_i)\|_2} - \frac{z_i}{\|z_i\|_2} \right \|_2 \\
& = \| \sum_{i=1}^n P_{z_i^\perp}(h)\|_2 - \|h\|_2 \sum_{i=1}^n \sqrt{2} \sqrt{1- \frac{\|P_{u^\perp} (z_i)\|_2}{\|z_i\|_2}}
\end{align*}
%The right-hand term can be handled similarly to above, since $v$ and $w_i$ are nearly orthogonal, so the resulting sum is $<< n \|h\|_2$. The left-hand term is exactly the same as occurs in the triangles inequality of the original ShapeFit paper, except with a different constant multiplying $t_i - p_i$. We follow the argument from Lemma ... in [cite ShapeFit] below:

Letting $X_i = \frac{\| P_{V_i^\perp}(t_i-p_i) \|_2}{\|t_i-p_i\|_2}$, $Y_i = \frac{\| P_{(x-y)^\perp} (x+y - t_i-p_i) \|_2}{\|x+y - t_i-p_i\|_2}$, and $Z_i = \left \|\sum_{i=1}^n (t_i-p_i)(t_i-p_i)^*\right\|_{\text{op}}$, we have shown that for any $h \perp x-y$,
\begin{align}
\sum_{i=1}^n \| P_{W_i^\perp}(h)\|_2  \notag
\geq \sum_{i=1}^n &\left \| P_{(x+y-t_i - p_i)^\perp}(h)\right \|_2 \\&-  \|h\|_2 \sum_{i=1}^n \sqrt{2}\left[ \sqrt{1 - X_i} + \sqrt{1 - Y_i} \right]  - \frac{1}{\min_i(\|t_i - p_i\|_2^2)}Z_i\|h\|_2 \label{pwi-perp}
\end{align}
We will separately bound the first term and last two terms with high probability.  



We now show that the first term of \eqref{pwi-perp} is bounded below by $0.3 n \|h\|_2$ with high probability.

Because $t_i + p_i =^d \sqrt{2}t_i$, it suffices to show that with high probability
\[
\left \| \sum_{i=1}^n P_{(x+y-\sqrt{2}t_i)^\perp}(h) \right\|_2 \geq   0.3 n\|h\|_2.
\]
Let $v = x+y$ and $w_i = -\sqrt{2}t_i$. Note that

\begin{align*}
\left \|  \sum_{i=1}^n P_{(v+w_i)^\perp}(h)\right\|_2 
&= \left \| \sum_{i=1}^n \left(h - \frac{1}{\|v+w_i\|_2^2}\langle h, v+ w_i \rangle (v+ w_i) \right)\right \|_2\\
&\geq n\|h\|_2 - \left \| \sum_{i=1}^n \frac{1}{\|v+w_i\|_2^2} (v+ w_i)(v+ w_i)^* h\right \|_2  \\
&\geq n\|h\|_2 - \left \| \sum_{i=1}^n \frac{1}{\|v+w_i\|_2^2} (v+ w_i)(v+ w_i)^* \right\|_{\text{op}} \|h\|_2\\
& \geq \|h\|_2 \left[n - \frac{1}{\min_{i}\|v+w_i\|_2^2} \left \| \sum_{i=1}^n (v+ w_i)(v+ w_i)^* \right\|_{\text{op}} \right],
\end{align*}
where in the last inequality we used 
\[
\sum_{i=1}^n \frac{1}{\|v+w_i\|_2^2} (v+ w_i)(v+ w_i)^* \preceq \frac{1}{\min_{i}\|v+w_i\|_2^2}  \sum_{i=1}^n (v+ w_i)(v+ w_i)^*.
\]
Now, let $A = \sum_{i=1}^n e_i w_i^* \in \mathbb{R}^{n\times d}$. We have  
\begin{align*}
\left \| \sum_{i=1}^n (v+ w_i)(v+ w_i)^* \right\|_{\text{op}} &= \left \| \sum_{i=1}^n (vv^* + vw_i^* + w_i v^* + w_i w_i^*) \right\|_{\text{op}}\\
& \leq n\|vv^*\|_{\text{op}} + \left\| v \left( \sum_{i=1}^n w_i \right)^* +  \left( \sum_{i=1}^n w_i \right)v^* \right\|_{\text{op}} + \left \| \sum_{i=1}^n w_i w_i^* \right\|_{\text{op}}\\
& \leq n \|v\|_2^2 + 2\|v\|_2\left \| \sum_{i=1}^n w_i \right\|_2 + \left \| \sum_{i=1}^n w_i w_i^* \right\|_{\text{op}}\\
& = n \|v\|_2^2 + 2\|v\|_2\left \| \sum_{i=1}^n w_i \right\|_2 + \sigma_{\max}(A)^2.
\end{align*}
Thus,
\[
\sum_{i=1}^n \left \| P_{(v+w_i)^\perp}(h)\right \|_2 \geq \|h\|_2\left[n - \frac{n \|v\|_2^2 + 2\|v\|_2\left \| \sum_{i=1}^n w_i \right\|_2 + \sigma_{\max}(A)^2}{\min_{i}\|v+w_i\|_2^2} \right].
\]

Now, consider the event
\[
E_1 = \left\{ \min_{i} \|v+w_i\|_2^2 \geq 4d\beta_1, \quad \|v\|_2^2 \leq 2d\beta_2, \quad \left \| \sum_{i=1}^n w_i \right\|_2^2 \leq 2nd\beta_3, \quad \sigma_{\max}(A)^2 \leq 2n\beta_4  \right\}
\]
On $E_1$ we have,
\begin{align*}
\sum_{i=1}^n \left \| P_{(v+w_i)^\perp}(h)\right \|_2 &\geq \|h\|_2 \left[ n - \frac{1}{4d\beta_1}\left( 2nd\beta_2 + 2\sqrt{2d\beta_2}\sqrt{2}\sqrt{nd}\sqrt{\beta_3} + 2n\beta_4  \right) \right]\\
&= \|h\|_2 \left[n - \frac{1}{2}n \frac{\beta_2}{\beta_1} - \frac{4d\sqrt{n} \sqrt{\beta_2\beta_3}}{4d\beta_1}-\frac{\beta_4 }{2d\beta_1}n \right]\\
& = \|h\|_2 \left[n\left(1 - \frac{1}{2} \frac{\beta_2}{\beta_1} - \frac{\beta_4}{2d\beta_1}- \frac{1}{\sqrt{n}}\frac{ \sqrt{\beta_2\beta_3}}{\beta_1} \right)\right]
\end{align*}
Now, let $\beta_1 = 1-\frac{1}{100}, \quad \beta_2 = \beta_3 = 1 + \frac{1}{100}, \quad \beta_4 = \frac{1}{5}d \beta_1$.  This  gives
\[
\frac{1}{2} \frac{\beta_2}{\beta_1} = 1/2 +1/99, \quad \frac{\beta_4}{2d\beta_1} = 1/10,  \quad \frac{1}{\sqrt{n}}\frac{  \sqrt{\beta_2\beta_3}}{\beta_1} < \frac{2}{\sqrt{n}}.
\]
Assuming $n \geq 550$, we see that on $E_1$,
\[\sum_{i=1}^n \left \| P_{(x+y-t_i - p_i)^\perp}(h)\right \|_2 \geq 0.3 n \|h\|_2. \]


Now, we bound $\P(E_1)$. 
Note that $ \frac{1}{4}\|v+w_i\|_2^2 =^d \frac{1}{2}\|v\|_2^2 =^d \frac{1}{2n} \left \| \sum_{i=1}^n w_i \right\|_2^2 =^d \chi^2(d)$ and $\frac{1}{\sqrt{2}}A$ is a random $n\times d$ matrix with i.i.d. $\mathcal{N}(0,1)$ entries. Thus, by applying Lemma \ref{lem:conc1}, we have 
\[
\P \Bigl( 4d(1-\epsilon) \leq \|v+w_i\|_2^2 \leq 4d(1+\epsilon) \Bigr) \geq 1 - e^{-c \epsilon^2 d}
\]

\[
\P\left( 2d(1-\epsilon) \leq \|v\|_2^2 \leq 2d(1+\epsilon) \right) \geq 1 - e^{-c \epsilon^2 d}
\]

\[
\P\left( 2nd(1-\epsilon) \leq \left \| \sum_{i=1}^n w_i \right\|_2^2 \leq 2nd(1+\epsilon) \right) \geq 1 - e^{-c \epsilon^2 d},
\]
where $c>0$ is a universal constant. Also by taking $t =2  \sqrt{d}$ in Lemma \ref{lem:conc2} we get
\[
\P \left( \sigma_{\text{max}}\Bigl(\frac{1}{\sqrt{2}} A\Bigr) \geq \sqrt{n} +3\sqrt{d} \right) \leq 2e^{-2d}
\]
We have
\[
\P  \left(\sigma_{\max} (\frac{1}{\sqrt{2}} A) \geq \sqrt{n \beta_4}\right) \leq \P\left( \sigma_{\max} (\frac{1}{\sqrt{2}} A) \geq \sqrt{n} + 3\sqrt{d} \right) \leq 2e^{-2d}
\]
whenever $\sqrt{n} + 3\sqrt{d} \leq \sqrt{n \beta_4}$, or equivalently $(\sqrt{\frac{\beta_1}{5}}\sqrt{d} -1) \sqrt{n} \geq 3 \sqrt{d}$,  %\left( \frac{2\sqrt{d}}{\sqrt{d/2} -1} \right)^2
which holds when $n \geq 550$ and $d \geq 10$. Thus for $n\geq 550$, we have
\[
\mathbb{P}(E_1) \geq 1- 2n e^{-cd} .
\]

We now show that the second term of \eqref{pwi-perp} is bounded above by $0.2 n \|h\|_2$ with high probability.
Define the event
\[
E_2 = \left\{ X_i \geq 1 - \frac{1}{800}, \quad Y_i \geq 1 - \frac{1}{800}, \quad \frac{1}{\min_i(\|t_i - p_i\|_2^2)}Z_i \leq 0.1 n, \quad i = 1,2\ldots n \right\}
\]
On $E_2$, we have
\[
 \|h\|_2 \sum_{i=1}^n \sqrt{2}\left[ \sqrt{1 - X_i} + \sqrt{1 - Y_i} \right]  + \frac{1}{\min_i(\|t_i - p_i\|_2^2)}Z_i\|h\|_2 \leq 0.2 n \|h\|_2.
\]
We now estimate $\mathbb{P}(E_2)$. For $X_i$, since $ (t_i - p_i)$ is independent from $(x-y, p_i + t_i - x-y)$, we can view the latter as fixed. That is, by conditioning on $V_i$, and applying a rotation $R$ such  that $R(V_i) = \Span(e_1,e_2)$, we have
\[
 \frac{\| P_{V_i^\perp} (t_i-p_i) \|_2}{\|t_i-p_i\|_2} =^d \sqrt{\frac{\sum_{j=1}^{d-2}t_i (j)^2}{\sum_{j=1}^{d}t_i (j)^2}} % \quad \approx \quad \sqrt{\frac{d-2}{d}} \quad \text{whp,}
\] 
where $t_i(j)$ is the $j$th entry of $t_i$.  As $\sum_{j=1}^{d-2}t_i (j)^2 \sim \chi^2_{d-2}$ and $\sum_{j=1}^{d}t_i (j)^2 \sim \chi^2_{d}$, Lemma \ref{lem:conc1} can be repeatedly applied to give $\P(X_i \geq 1 - \frac{1}{800} \text{ for all } i) \geq 1 - 2 n e^{-cd}$. 
%\blue{(why is this here?) which gives
%\[
%2 \|h\|_2 \sum_{i=1}^n \sqrt{2}\sqrt{1 - \frac{\| P_{V_i^\perp} (t_i-p_i) \|_2}{\|t_i-p_i\|_2}} \quad \approx \quad  \frac{2\sqrt{2}}{\sqrt{d}} n\|h\|_2.
%\]
%}
A similar argument gives $\P(Y_i \geq 1 - \frac{1}{800} \text{ for all } i) \geq 1 - 2n e^{-cd}$ because  $x-y$ and $x+y-(t_i+p_i)$ are independent.


We now bound the probability of the third condition in the definition of $E_2$.  Note that 
\[
\frac{1}{\min_i(\|t_i - p_i\|_2^2)}\left \|\sum_{i=1}^n (t_i-p_i)(t_i-p_i)^*\right\|_{\text{op}} =^d \frac{1}{\min_i(\|t_i \|_2^2)}\left \|\sum_{i=1}^n t_i t_i^*\right\|_{\text{op}}
\]
Let $B = \sum_{i=1}^n e_i t_i^*$.  By Lemma \ref{lem:conc2}, $\| \sum_{i=1}^n t_i t_i^*\|_\text{op} = \sigma_\text{max}(B)^2 \geq n \Bigl(1 + 3 \sqrt{\frac{d}{n}} \Bigr)^2$ with probability at least $1 - 2 e^{-2d}$.
By Lemma \ref{lem:conc1}, $\|t_i\|_2^2 \geq d(1-\eps)$ for all $i$ with probability at least $1 - n e^{-c \eps^2 d}$.  We conclude
\[
 \frac{1}{\min_i(\|t_i \|_2^2)}\left \|\sum_{i=1}^n t_i t_i^*\right\|_{\text{op}} \leq \frac{n \Bigl(1 + 3 \sqrt{\frac{d}{n}} \Bigr)^2}{d (1-\eps)}
\]
with probability at least $1 - 2n e^{-c \eps^2 d}$.
If $\eps = 0.01, d \geq 40, n \geq 10d$, we have
\[
\P\Bigl( \frac{1}{\min_i(\|t_i - p_i\|_2^2)}Z_i \leq 0.1 n \Bigr)  \geq 1 - 2n e^{-cd}
\]
Hence, if $d \geq 40, n \geq 10 d$,
\[
\P(E_2) \geq 1 - 6n e^{-cd}.
\]


%\blue{MORE HERE}
%For the \red{third condition in the definition of $E_2$}, we have with high probability 
%\[
%\frac{1}{\min_i(\|t_i - p_i\|_2^2)}\left \|\sum_{i=1}^n (t_i-p_i)(t_i-p_i)^*\right\|_{\text{op}} \|h\|_2 \quad \lesssim \quad \frac{n}{d} \left( 1 + \frac{18d}{n} + \frac{12\sqrt{d}}{\sqrt{n}} \right) \|h\|_2
%\]
%for $n \gg d$ and $n,d$ large enough. Overall, by taking $\epsilon$ small enough in repeated applications of Lemma \ref{lem:conc1}, and by applying Lemma \ref{lem:conc2} we have that $$%\mathbb{P}(E_2) \geq 1 - 4n e^{-cd} - 2e^{-d}.$$ 
In conclusion, there exist positive integers $d_0$ and $n_0$ such that for all $d \geq d_0$, $n \geq n_0$, $n \geq 10 d$, and all $h \perp x-y$,
\[
\mathbb{P}\left(\left \| \sum_{i=1}^n P_{{W_i}^\perp} (h) \right\|_2 \geq \frac{1}{10}\|h\|_2 \right) \geq 1- \mathbb{P}\left[(E_1 \cap E_2)^c\right] \geq 1- 6ne^{-cd}
\]
for some $c > 0$, which implies the statement of the lemma. 



\end{proof}


%%
%%
%%
%%	SECTION : main
%%
%%
%%

%\begin{theorem} \label{thm:mainthm}
%For all $p,\beta,c_{0}$ and $c_{1}$, there exists
%$\varepsilon$ such that under the assumptions stated above, we have
%$R(T)>R(\Tnot)$ for all $T\neq \Tnot$ such that $L(T)=1$.
%\end{theorem}

\subsection{Random graphs are $p$-typical with high probability} \label{sec-random-graph}

We prove that Condition 1 of Theorem~\ref{thm:mainthm} holds with high probability.

\begin{lemma} \label{lem:ptyp}
There exists an absolute constant $c > 0$ such that for all positive real numbers $p \le 1$ satisfying $n_2 p \ge 2\log (e n_1)$ and $n_1 p \ge 2\log (e n_2)$, $G(n_1,n_2;p)$ is $p$-typical with probability at least $1 - n_1 n_2 2^{n_1 + n_2} e^{-pn_1n_2/4} - n_1^2 n_2 e^{-\Omega(n_2p^2)} - n_1 n_2^2 e^{-\Omega(n_1p^2)}$.
\end{lemma}
\begin{proof}
Let $V_1$ and $V_2$ be vertex sets of sizes $|V_1| = n_1$ and $|V_2| = n_2$. 
Throughout the proof, we let $V_1 \cup V_2$ be the bipartition of the random graph $G(n_1,n_2;p)$.
The bipartite graph $G(n_1,n_2;p)$ is
not connected only if there exist partitions $V_1 = V_{1,1} \cup V_{1,2}$ 
and $V_2 = V_{2,1} \cup V_{2,2}$ such that the sets $V_{1,1} \cup V_{2,1}$
and $V_{1,2} \cup V_{2,2}$ are both non-empty and have no edges between them. 
Let $|V_{1,1}| = k_1, |V_{2,1}| = k_2$, $|V_{1,2}| = n_1 - k_1$ and $|V_{2,2}| = n_2 - k_2$. 
For fixed $k_1, k_2$, by the union bound, 
the probability that there exists a  partition as above is at most
\begin{align} \label{eq:disconnected}
	{n_1 \choose k_1} {n_2 \choose k_2} (1-p)^{k_1 (n_2 - k_2) + k_2 (n_1 - k_1)}.
\end{align}
If $k_1 \le \frac{n_1}{2}$ and $k_2 \le \frac{n_2}{2}$, then by Stirling's formula, \eqref{eq:disconnected}
is at most
\[
	\left(\frac{en_1}{k_1}\right)^{k_1} 
	\left(\frac{en_2}{k_2}\right)^{k_2}
	(1-p)^{(k_1 n_2 + k_2 n_1)/2}
	\le
	\left(\frac{en_1}{k_1} e^{-n_2p/2}\right)^{k_1} 
	\left(\frac{en_2}{k_2} e^{-n_1p/2}\right)^{k_2}.
\]
If $k_1 > \frac{n_1}{2}$ and $k_2 > \frac{n_2}{2}$, then 
let $\ell_1 = n_1 - k_1$ and $\ell_2 = n_2 - k_2$. Then \eqref{eq:disconnected} is at most
\[
	\left(\frac{en_1}{\ell_1}\right)^{\ell_1} 
	\left(\frac{en_2}{\ell_2}\right)^{\ell_2}
	(1-p)^{(\ell_1 n_2 + \ell_2 n_1)/2}
	\le
	\left(\frac{en_1}{\ell_1} e^{-n_2p/2}\right)^{\ell_1} 
	\left(\frac{en_2}{\ell_2} e^{-n_1p/2}\right)^{\ell_2}.
\]
If ($k_1 \le \frac{n_1}{2}$ and $k_2 > \frac{n_2}{2}$)
or ($k_1 > \frac{n_1}{2}$ and $k_2 \le \frac{n_2}{2}$), then, by ${n \choose k} \leq 2^n$ for all $0 \leq k \leq n$, \eqref{eq:disconnected}
is at most
\[
	2^{n_1 + n_2} (1-p)^{n_1 n_2/4}
	\le 2^{n_1 + n_2} e^{-pn_1n_2/4}.
\]

Hence the probability that $G(n_1, n_2;p)$ is disconnected is at most
\begin{align*}
	&\, \sum_{k_1=1}^{\lfloor n_1/2 \rfloor}
	\sum_{k_2=0}^{\lfloor n_2/2 \rfloor}
	\left(\frac{en_1}{k_1} e^{-n_2p/2}\right)^{k_1} 
	\left(\frac{en_2}{k_2} e^{-n_1p/2}\right)^{k_2}
	+
	\sum_{k_1=0}^{\lfloor n_1/2 \rfloor}
	\sum_{k_2=1}^{\lfloor n_2/2 \rfloor}
	\left(\frac{en_1}{k_1} e^{-n_2p/2}\right)^{k_1} 
	\left(\frac{en_2}{k_2} e^{-n_1p/2}\right)^{k_2}
	\\
	&\,+ \sum_{\ell_1=1}^{\lfloor n_1/2 \rfloor}
	\sum_{\ell_2=0}^{\lfloor n_2/2 \rfloor}
	\left(\frac{en_1}{\ell_1} e^{-n_2p/2}\right)^{\ell_1} 
	\left(\frac{en_2}{\ell_2} e^{-n_1p/2}\right)^{\ell_2}
	+
	\sum_{\ell_1=0}^{\lfloor n_1/2 \rfloor}
	\sum_{\ell_2=1}^{\lfloor n_2/2 \rfloor}
	\left(\frac{en_1}{\ell_1} e^{-n_2p/2}\right)^{\ell_1} 
	\left(\frac{en_2}{\ell_2} e^{-n_1p/2}\right)^{\ell_2}
	\\
	&\,+ n_1 n_2 2^{n_1 + n_2} e^{-pn_1n_2/4},
\end{align*}
where the indeterminate factors in the sums corresponding to $k_1=0$, $k_2=0$, $\ell_1=0$, or $\ell_2=0$ are taken to be unity.
Since $n_2 p \ge 2\log (e n_1)$ and $n_1 p \ge 2\log (e n_2)$, the four sums above are maximized
at $(k_1, k_2)=(1,0), (0,1)$, $(\ell_1, \ell_2)=(1,0), (0,1)$, respectively. Therefore
the probability that $G(n_1, n_2;p)$ is disconnected is at most
\begin{align*}
	&\, 
	2n_1 n_2
	\cdot en_1 \cdot e^{-n_2p/2}
	+
	2n_1 n_2
	\cdot en_2 \cdot e^{-n_1p/2}
	+ n_1 n_2 2^{n_1 + n_2} e^{-pn_1n_2/4}. 
\end{align*}

For a fixed vertex $v \in V_1$, the expected value of $\deg(v)$ is $n_2 p$, 
and for a pair of vertices $v,w \in V_1$, the expected value of
the codegree of $v$ and $w$ is $n_2 p^{2}$. Therefore by 
Chernoff's inequality (see Fact 4 from  \cite{AngluinValiant}) and a union bound, 
the probability that all vertices in $V_1$ have degree between $\frac{1}{2}n_2 p$ and $2n_2p$, 
and all pairs of vertices in $V_1$ have codegree between $\frac{1}{2} n_2p^2$ and $2n_2p^2$ is
$1 - n_1^2 e^{-\Omega(n_2p^2)}$.
Similarly, the probability that all vertices in $V_2$ have degree between $\frac{1}{2}n_1 p$ and $2n_1p$, 
and all pairs of vertices in $V_2$ have codegree between $\frac{1}{2} n_1p^2$ and $2n_1p^2$ is
$1 - n_2^2 e^{-\Omega(n_1p^2)}$. The conclusion follows by taking a union bound over
all events.
\end{proof}





\subsection{Proof of Theorem \ref{thm:bipartite-random}} \label{sec-proof-random-theorem}
We can now prove the high dimensional recovery theorem, which we state here again for convenience: %\red{$K_0$ is a constant from Vlad's lemma}
\addtocounter{theorem}{-2}
\begin{theorem} 
Let $N = \max(\nl,\ns), n = \min(\nl, \ns)$. Let $G(\Vl \cup \Vs,E)$ be drawn from a bipartite-\erdosrenyi graph with $p >0$. Take $ \tnot_1, \ldots \tnot_{\nl}, \pnot_1, \ldots, \pnot_{\ns} \sim \mathcal{N}(0, I_{d \times d})$ to be independent from each other and $G$. Then, there exist absolute constants $c,c_3, C>0$ such that for $\gamma = c_3 p^4$, if \\
$$\max\left(\frac{1}{c_3 p^4}, Cd, \frac{2 \log(eN)}{p}, \Omega(c_3 \log^2 N)\right) \leq  n \le N \leq e^{\frac{1}{8}c d }$$ 
and $d = \Omega(1)$, then there exists an event with probability at least  $1 - O( e^{-\Omega(\frac{1}{2}c_3^{-1/2} n^{1/2})} + e^{-\frac{1}{2} cd})$, on which the following holds:\\[1em]
For all subgraphs $\Eb$ satisfying $\max_{i\in [\nl]} \deg_b(\ti) \leq \gamma \ns$ and $\max_{j \in [\ns]} \deg_b(\pj) \leq \gamma \nl$ and all pairwise direction corruptions  $\vij \in \mathbb{S}^{d-1}$ for $ij \in \Eb$,  the convex program \eqref{shapefit} has a unique minimizer equal to $\left\{\alpha \{\toi -\tpbar \}_{i \in [\nl]}, \alpha \{\poi -\tpbar\}_{j \in [\ns]}  \right\}$ for some positive $\alpha$ and for $\tpbar = \frac{1}{\nl + \ns} \left(\sum_{i \in [\nl]} \toi + \sum_{j \in [\ns]} \poj \right)$. 
\end{theorem}
\begin{proof}
Let $c$ be minimum of the constants from Lemmas \ref{lem-conds34} and \ref{lem-c4-well-distributed}. Let $K_0$ be the constant from Lemma \ref{lem-c4-well-distributed}.
It is enough to verify that $G$, $T$ and $E_b$ in the assumption of the present theorem satisfy the deterministic conditions 1--6 in Theorem \ref{thm:mainthm}, with appropriate constants $p, \beta, c_0, \epsilon, c_1$, and with the purported probability.
By Lemma \ref{lem:ptyp}, Condition 1 holds with probability at least 
$$1 - n_l n_s 2^{n_l + n_s} e^{-pn_ln_s/4} - n_l^2 n_s e^{-\Omega(n_sp^2)} - n_l n_s^2 e^{-\Omega(n_lp^2)} = 1 - O(N^3 e^{-\Omega(np^2)})$$
 if $np \geq 2 \log(eN)$.  Condition 2 holds with probability $1$. By Lemma \ref{lem-conds34}, Condition 3 holds for $c_0 = \frac{9}{10}$ with probability at least $1 - 2 \nl \ns e^{-cd}$, and Condition 4 holds for $\beta = \frac{1}{4}$ with probability at least $1 - 22 \nl^2 \ns^2 e^{-c d}$. 
By Lemma \ref{lem-c4-well-distributed}, Condition 5 holds for $c_1 = \frac{1}{20}$ with probability $1 - O( \nl^2 \ns^2 e^{-cd})$ if $n \geq \max(K_0, 160d) $, and $d \geq d_0$.
 Thus, Conditions 1--5 hold together with probability at least
\[
1 -  O(N^4 e^{-cd} +N^3 e^{-\Omega(n p^2)}).
\]
%where $c>0$ is an absolute constant.

Take $\gamma = c_3 p^4 \leq \frac{p^4}{10^{11}}$.  Because $\gamma \leq \frac{\beta c_0 c_1^2 p^4}{384\cdot204\cdot64}$, Theorem \ref{thm:mainthm} implies that recovery via ShapeFit is guaranteed.  Note that the conditions $\max_{i\in\Vl} \deg_b(i) \leq \gamma \ns$ and $\max_{j \in \Vs} \deg_b(j) \leq \gamma \nl$ are nontrivial when $p \geq c_3^{-1/4} n^{-1/4}$.  Using this inequality, we have $N^3 e^{-\Omega(np^2)} \leq N^3 e^{-\Omega(c_3^{-1/2} n^{1/2})} \leq e^{-\Omega(\frac{1}{2} c_3^{-1/2} n^{1/2})}$ if $n = \Omega(c_3 \log^2 N) $ and $N^4 e^{-cd} \le e^{-cd/2}$ if $N \leq e^{\frac{1}{8}c d}$. Thus, the probability of exact recovery via ShapeFit, uniformly in $E_b$ and $v_{ij}$ satisfying the assumptions of the theorem, is at least
\[
1- O(e^{-\Omega(\frac{1}{2}c_3^{-1/2} n^{1/2})} + e^{-\frac{1}{2}c d}). \qedhere
\]

%probability at least, that (1) holds with probability at least... Then the conclusion of the theorem guarantees that taking $\gamma < \frac{\beta c_0 c_1^2 p^4}{10^6}$, guarantees recovery via ShapeFit. 


\end{proof}


































































































